\documentclass[11pt]{article}
\usepackage{mathblatt}
\begin{document}
\blattkopf*{Prüfungsheft Abitur GK · Analytische Geometrie}{Prüfungsheft}{Prüfungsheft Abitur GK · Analytische Geometrie · Prüfungsheft · PUN-H1}
\verzeichniszeile{\verz{t1}{Punkte und Strecken im Koordinatensystem (Nr.~1–18)} \verztrenn \verz{t2}{Vektoren und Rechenoperationen (Nr.~19–25)} \verztrenn \verz{t3}{Geraden (Nr.~26–36)} \verztrenn \verz{t4}{Ebenen (Nr.~37–48)} \verztrenn \verz{t5}{Lagebeziehungen (Nr.~49–56)} \verztrenn \verz{t6}{Schnittmengen (Nr.~57–64)} \verztrenn \verz{t7}{Abstände (Nr.~65–72)} \verztrenn \verz{t8}{Skalarprodukt und Winkel (Nr.~73–80)} \verztrenn \verz{t9}{Orthogonalität (Nr.~81–88)} \verztrenn \verz{t10}{Flächeninhalt und Volumen im Raum (Nr.~89–99)}}
% Punkte und Strecken im Koordinatensystem – bank/punkte-und-strecken-im-koordinatensystem/ (zusammenbau v0.4, Heft)
\einheitenkopf[t1]{Punkte und Strecken im Koordinatensystem}
\setcounter{aufgabe}{0}

% Anlauf (ohne Prüfkennung)
\begin{aufgabe}{Darstellen und Lage lesen – Anlauf}
\begin{teile}
\teil $P(4 | -1 | 0)$ – was sagt die Null? \\ \kreuz{$P$ liegt auf der $x_3$-Achse} \\ \kreuz{$P$ liegt in der $x_1x_2$-Ebene} \\ \kreuz{$P$ liegt in der $x_2x_3$-Ebene}
\teil $A(5 | 3 | 1)$ und $B(2 | 1 | 4)$ eintragen, $C$ ablesen.

\begin{ksys3} \rpunkt{4}{5}{2}{C} \end{ksys3}
C( \leerfeld | \leerfeld | \leerfeld )
\teil Zwei Masten stehen senkrecht in $F_1(-2 | -2 | 0)$ ($4$ m hoch) und $F_2(5 | 3 | 0)$ ($3$ m hoch); ein Seil verbindet die Spitzen (1 LE = 1 m). Zeichne Masten und Seil.

\begin{ksys3} \end{ksys3}
\end{teile}
\end{aufgabe}

\begin{aufgabe}{Darstellen und Lage lesen – Prüfungsaufgaben}
\begin{teile}
\teil Ein Holzwürfel mit Kante $4$ hat eine Ecke im Ursprung und liegt im ersten Oktanten (alle Koordinaten $\ge 0$). Zeichne das Viereck $P(4 | 0 | 1)$, $Q(1 | 4 | 0)$, $R(0 | 4 | 2)$, $S(3 | 0 | 4)$ ein. \hfill (Abitur 2019 GK)

\begin{ksys3} \rquader{0,0,0}{4}{4}{4} \end{ksys3}
\teil Ein Glaswürfel mit Kante $3$ hat eine Ecke im Ursprung und liegt im ersten Oktanten (alle Koordinaten $\ge 0$). Zeichne das Viereck $P(3 | 1 | 0)$, $Q(3 | 3 | 2)$, $R(1 | 3 | 3)$, $S(0 | 0 | 2)$ ein. \hfill (Abitur 2019 GK)

\begin{ksys3} \rquader{0,0,0}{3}{3}{3} \end{ksys3}
\teil Ein Würfel mit Kante $4$ hat eine Ecke im Ursprung und liegt im ersten Oktanten (alle Koordinaten $\ge 0$). Zeichne das Viereck $P(0 | 3 | 0)$, $Q(4 | 4 | 1)$, $R(4 | 1 | 4)$, $S(1 | 0 | 4)$ ein. \hfill (Abitur 2019 GK)

\begin{ksys3} \rquader{0,0,0}{4}{4}{4} \end{ksys3}
\teil Ein Quader hat eine Ecke im Ursprung $O$ und die Kanten $5$, $6$ und $3$ entlang der $x_1$-, $x_2$- und $x_3$-Achse. $P(5 | 2 | 3)$ liegt auf einer Kante. Zeichne das Dreieck mit den Ecken $O$, $P$ und $R(0 | 6 | 3)$ ein. \hfill (Abitur 2026 GK)

\begin{ksys3} \rquader{0,0,0}{5}{6}{3} \end{ksys3}
\teil Ein Quader hat eine Ecke im Ursprung $O$ und die Kanten $6$, $3$ und $4$ entlang der $x_1$-, $x_2$- und $x_3$-Achse. $P(6 | 0 | 1)$ liegt auf einer Kante. Zeichne das Dreieck mit den Ecken $O$, $P$ und $T(0 | 3 | 4)$ ein. \hfill (Abitur 2026 GK)

\begin{ksys3} \rquader{0,0,0}{6}{3}{4} \end{ksys3}
\teil Ein Quader hat eine Ecke im Ursprung $O$ und die Kanten $3$, $6$ und $3$ entlang der $x_1$-, $x_2$- und $x_3$-Achse. $P(3 | 2 | 3)$ liegt auf einer Kante. Zeichne das Dreieck mit den Ecken $O$, $P$ und $T(3 | 6 | 0)$ ein. \hfill (Abitur 2026 GK)

\begin{ksys3} \rquader{0,0,0}{3}{6}{3} \end{ksys3}
\end{teile}
\end{aufgabe}

\begin{aufgabe}{Darstellen und Lage lesen – Prüfungsaufgaben – weiter}
\begin{teile}
\teil Begründe: $P(-3 | 4 | 3)$ und $Q(5 | -1 | 2)$ liegen auf derselben Seite der $x_1x_2$-Ebene. \hfill (Abitur 2025 GK)
\teil Begründe: Das Dreieck $A(6 | 0 | 1)$, $B(6 | 3 | 5)$, $C(6 | -3 | 5)$ liegt parallel zur $x_2x_3$-Ebene und symmetrisch zur $x_1x_3$-Ebene. \hfill (Abitur 2026 GK)
\teil Begründe: $R(1 | -2 | -3)$ und $T(4 | 5 | 2)$ liegen auf verschiedenen Seiten der $x_1x_2$-Ebene. \hfill (Abitur 2025 GK)
\teil Der Grundriss eines Pavillons ist ein Achteck in der $x_1x_2$-Ebene, symmetrisch zur $x_1$- und zur $x_2$-Achse. Gegeben sind $A(6 | 0)$, $B(5 | 3)$, $C(0 | 4)$. Zeichne das vollständige Achteck. \hfill (Abitur 2025 GK)

\begin{ksys}[xmin=-7,xmax=7,ymin=-7,ymax=7,xlabel=$x_1$,ylabel=$x_2$] \punkt{6}{0}{A} \punkt{5}{3}{B} \punkt{0}{4}{C} \end{ksys}
\teil Der Beet ist ein Achteck in der $x_1x_2$-Ebene, symmetrisch zur $x_1$- und zur $x_2$-Achse. Gegeben sind $A(4 | 0)$, $B(3 | 2)$, $C(0 | 3)$. Zeichne das vollständige Achteck. \hfill (Abitur 2025 GK)

\begin{ksys}[xmin=-7,xmax=7,ymin=-7,ymax=7,xlabel=$x_1$,ylabel=$x_2$] \punkt{4}{0}{A} \punkt{3}{2}{B} \punkt{0}{3}{C} \end{ksys}
\teil Der Bodenplatte ist ein Achteck in der $x_1x_2$-Ebene, symmetrisch zur $x_1$- und zur $x_2$-Achse. Gegeben sind $A(6 | 0)$, $B(2 | 5)$, $C(0 | 6)$. Zeichne das vollständige Achteck. \hfill (Abitur 2025 GK)

\begin{ksys}[xmin=-7,xmax=7,ymin=-7,ymax=7,xlabel=$x_1$,ylabel=$x_2$] \punkt{6}{0}{A} \punkt{2}{5}{B} \punkt{0}{6}{C} \end{ksys}
\end{teile}
\end{aufgabe}

\begin{aufgabe}{Darstellen und Lage lesen – Prüfungsaufgaben – weiter}
\begin{teile}
\teil Eine Pyramide hat den quadratischen Boden $ABCD$ mit der Seite $3$ cm; die Spitze $S$ liegt $4$ cm senkrecht über $A$. Im Gitter sind der Boden und die Spitze des Seitendreiecks $ABS$ markiert. Zeichne die drei übrigen Seitendreiecke ein. \hfill (Abitur 2024 GK)

\begin{ksys}[xmin=-5,xmax=9,ymin=-5,ymax=9] \punkt{0}{0}{A} \punkt{3}{0}{B} \punkt{3}{3}{C} \punkt{0}{3}{D} \punkt{0}{-4}{S} \end{ksys}
\teil Eine Pyramide hat den quadratischen Boden $ABCD$ mit der Seite $4$ cm; die Spitze $S$ liegt $3$ cm senkrecht über $A$. Im Gitter sind der Boden und die Spitze des Seitendreiecks $ABS$ markiert. Zeichne die drei übrigen Seitendreiecke ein. \hfill (Abitur 2024 GK)

\begin{ksys}[xmin=-4,xmax=10,ymin=-4,ymax=10] \punkt{0}{0}{A} \punkt{4}{0}{B} \punkt{4}{4}{C} \punkt{0}{4}{D} \punkt{0}{-3}{S} \end{ksys}
\teil Eine Pyramide hat den quadratischen Boden $ABCD$ mit der Seite $2$ cm; die Spitze $S$ liegt $1{,}5$ cm senkrecht über $A$. Im Gitter sind der Boden und die Spitze des Seitendreiecks $ABS$ markiert. Zeichne die drei übrigen Seitendreiecke ein. \hfill (Abitur 2024 GK)

\begin{ksys}[xmin=-2.5,xmax=5.5,ymin=-2.5,ymax=5.5] \punkt{0}{0}{A} \punkt{2}{0}{B} \punkt{2}{2}{C} \punkt{0}{2}{D} \punkt{0}{-1.5}{S} \end{ksys}
\teil Die Pyramide $ABCS$ hat die Grundfläche $A(4 | 1 | 2)$, $B(1 | 5 | 2)$, $C(-1 | -1 | 2)$ und die Spitze $S(4 | -1 | 6)$. Eingezeichnet sind die Projektionen $A'$ und $B'$ in die $x_1x_2$-Ebene. Ergänze $C'$ und $S'$ und entscheide: Liegt der Höhenfußpunkt innerhalb der Grundfläche? \hfill (Abitur 2018 GK)

\begin{ksys3} \rpunkt*{4}{1}{0}{A'} \rpunkt*{1}{5}{0}{B'} \end{ksys3}
\teil Die Pyramide $ABCS$ hat die Grundfläche $A(5 | -1 | 1)$, $B(1 | 5 | 1)$, $C(-1 | -2 | 1)$ und die Spitze $S(1 | 1 | 6)$. Eingezeichnet sind die Projektionen $A'$ und $B'$ in die $x_1x_2$-Ebene. Ergänze $C'$ und $S'$ und entscheide: Liegt der Höhenfußpunkt innerhalb der Grundfläche? \hfill (Abitur 2018 GK)

\begin{ksys3} \rpunkt*{5}{-1}{0}{A'} \rpunkt*{1}{5}{0}{B'} \end{ksys3}
\end{teile}
\end{aufgabe}

% Anlauf (ohne Prüfkennung)
\begin{aufgabe}{Streckenlänge, Mittelpunkt, Teilpunkte – Anlauf}
\begin{teile}
\teil Wo muss der Haken an der Stange zwischen $A$ und $B$ hängen, damit er genau in der Mitte ist? Was ist gefragt? \\ \kreuz{eine Länge} \\ \kreuz{ein Punkt} \\ \kreuz{eine Figur}
\teil $A(2 | 1 | 3)$, $B(4 | 4 | 9)$: Länge von $AB$? \leerfeld LE
\teil Ein Zeltmast endet in $S(2 | 2 | 6)$; vier gleich lange Leinen führen zu Heringen am Boden, einer davon steckt in $H(4 | 5 | 0)$ (1 LE = 1 m). Je Leine kommen $40$ cm für die Knoten dazu. Wie viel Leine braucht man insgesamt? \leerfeld[m]
\end{teile}
\end{aufgabe}

\begin{aufgabe}{Streckenlänge, Mittelpunkt, Teilpunkte – Prüfungsaufgaben}
\begin{teile}
\teil Ein Gummiseil ist von $A(3 | 1 | 2)$ nach $B(9 | 9 | 2)$ gespannt (1 LE = 1 m) und dabei um $25\,\%$ länger als ungespannt. Wie lang ist es ungespannt? \hfill (Abitur 2020 GK) \\ \leerfeld[m]
\teil Ein Seil hängt belastet von $A(1 | 1 | 10)$ bis $B(3 | 7 | 1)$ (1 LE = 1 m) und ist dabei $10\,\%$ länger als unbelastet. Unbelastete Länge? \hfill (Abitur 2020 GK) \\ \leerfeld[m]
\teil Eine Feder zwischen $A(2 | 3 | 1)$ und $B(3 | 7 | 9)$ (1 LE = 1 cm) ist $50\,\%$ länger als in Ruhe. Länge in Ruhe? \hfill (Abitur 2020 GK) \\ \leerfeld[cm]
\teil Eine Pyramide ist symmetrisch zur $x_1x_3$-Ebene; $D(1 | 3 | 0)$ ist eine Grundecke, $S(3 | 0 | 6)$ die Spitze. Zeige: $W(2 | 1{,}5 | 3)$ ist der Mittelpunkt der Kante $DS$. Gib den Mittelpunkt $T$ der zu $DS$ symmetrischen Kante an. \hfill (Abitur 2026 GK) \\ T( \leerfeld | \leerfeld | \leerfeld )
\teil Eine Dachgaube ist symmetrisch zur $x_2x_3$-Ebene; $D(6 | 1 | 0)$ ist eine Ecke, $S(0 | 5 | 8)$ die Firstspitze. Zeige: $W(3 | 3 | 4)$ ist der Mittelpunkt der Kante $DS$. Gib den Mittelpunkt $T$ der zu $DS$ symmetrischen Kante an. \hfill (Abitur 2026 GK) \\ T( \leerfeld | \leerfeld | \leerfeld )
\teil Ein Kristall ist eine Doppelpyramide, symmetrisch zur $x_1x_2$-Ebene; $D(7 | 4 | 0)$ ist eine Ecke, $S(1 | 2 | 4)$ die obere Spitze. Zeige: $W(4 | 3 | 2)$ ist der Mittelpunkt der Kante $DS$. Gib den Mittelpunkt $T$ der zu $DS$ symmetrischen Kante an. \hfill (Abitur 2026 GK) \\ T( \leerfeld | \leerfeld | \leerfeld )
\teil Eine Strebe führt von $P(9 | 0 | 3)$ nach $Q(0 | 3 | 0)$; zwei Schellen teilen sie in drei gleich lange Stücke. Koordinaten der Schellen? \hfill (Abitur 2024 GK)
\teil Ein Kletterseil ist von Haken $A(1 | 4 | 2)$ zu Haken $C(11 | -1 | 7)$ gespannt. Der Karabiner $G$ teilt es so, dass $|AG| : |GC| = 2 : 3$. Koordinaten von $G$? \hfill (Abitur 2023 GK) \\ G( \leerfeld | \leerfeld | \leerfeld )
\teil Eine Stange reicht von $A(2 | 0 | 8)$ bis $C(6 | 8 | 0)$. Eine Leuchte $G$ sitzt so, dass $|AG| : |GC| = 3 : 1$. Koordinaten von $G$? \hfill (Abitur 2023 GK) \\ G( \leerfeld | \leerfeld | \leerfeld )
\teil $A(3 | 5 | 2)$ und $B(5 | 6 | 4)$ liegen auf einer Geraden $g$. Gib einen Punkt $D \ne B$ auf $g$ an, der von $A$ so weit entfernt ist wie $B$. \hfill (Abitur 2020 GK) \\ D( \leerfeld | \leerfeld | \leerfeld )
\teil $A(1 | -2 | 4)$ und $B(4 | 2 | 4)$ liegen auf einer Geraden $g$. Gib einen Punkt $D \ne B$ auf $g$ an, der von $A$ so weit entfernt ist wie $B$. \hfill (Abitur 2020 GK) \\ D( \leerfeld | \leerfeld | \leerfeld )
\teil $A(0 | 3 | 5)$ und $B(2 | 1 | 6)$ liegen auf einer Geraden $g$. Gib einen Punkt $D \ne B$ auf $g$ an, der von $A$ so weit entfernt ist wie $B$. \hfill (Abitur 2020 GK) \\ D( \leerfeld | \leerfeld | \leerfeld )
\end{teile}
\end{aufgabe}

\begin{aufgabe}{Streckenlänge, Mittelpunkt, Teilpunkte – Prüfungsaufgaben – weiter}
\begin{teile}
\teil $A(2 | -1 | 2)$; der Punkt $C$ ist festgelegt durch $\overrightarrow{OC} = 3 \cdot \overrightarrow{OA}$. Länge der Strecke $AC$? \hfill (Abitur 2018 GK) \\ \leerfeld
\teil $A(2 | 3 | 6)$; der Punkt $C$ ist festgelegt durch $\overrightarrow{OC} = - \overrightarrow{OA}$. Länge der Strecke $AC$? \hfill (Abitur 2018 GK) \\ \leerfeld
\teil $A(4 | 4 | 2)$; der Punkt $C$ ist festgelegt durch $\overrightarrow{OC} = \frac{1}{2} \cdot \overrightarrow{OA}$. Länge der Strecke $AC$? \hfill (Abitur 2018 GK) \\ \leerfeld
\teil Ein Tunnel verläuft geradlinig von $A(10 | 30 | 3)$ nach $B(40 | 70 | 3)$ (1 LE = 10 m). Ein Zug fährt ab $A$ mit $60$ km/h, gleichzeitig ein zweiter ab $B$ mit $90$ km/h. Welchen Weg im Tunnel legt der Zug ab $B$ zurück, bis sich beide begegnen? \hfill (Abitur 2019 GK) \\ \leerfeld[m]
\teil Ein Radweg führt geradlinig von $P(2 | 1 | 0)$ nach $Q(8 | 9 | 0)$ (1 LE = 100 m). Anna startet in $P$ mit $15$ km/h, Ben gleichzeitig in $Q$ mit $10$ km/h. Wie weit ist Anna gefahren, wenn sie sich treffen, und wo ist das? \hfill (Abitur 2019 GK)
\end{teile}
\end{aufgabe}

% Anlauf (ohne Prüfkennung)
\begin{aufgabe}{Dreiecke nachweisen – Anlauf}
\begin{teile}
\teil Parallelogramm $ABCD$: Welcher Vektor ist gleich $\overrightarrow{AB}$? \\ \kreuz{$\overrightarrow{CD}$} \\ \kreuz{$\overrightarrow{DC}$} \\ \kreuz{$\overrightarrow{BC}$}
\end{teile}
\end{aufgabe}

\begin{aufgabe}{Dreiecke nachweisen – Prüfungsaufgaben}
\begin{teile}
\teil Zeige: Das Dreieck $A(2 | 0 | 1)$, $B(6 | 2 | 1)$, $C(4 | 1 | 5)$ ist gleichschenklig. Nenne Basis und Schenkel. \hfill (Abitur 2023 GK)
\teil Zeige: Das Dreieck $A(0 | 2 | 3)$, $B(4 | 2 | 1)$, $C(2 | 5 | 2)$ ist gleichschenklig. Nenne Basis und Schenkel. \hfill (Abitur 2023 GK)
\teil Zeige: Das Dreieck $P(1 | 4 | 0)$, $Q(5 | 2 | 0)$, $R(5 | 7 | 2)$ ist gleichschenklig. Nenne Basis und Schenkel. \hfill (Abitur 2023 GK)
\teil Zeige: Das Dreieck $A(2 | 1 | 4)$, $B(4 | 5 | 2)$, $C(5 | 3 | 5)$ ist gleichschenklig. Nenne Basis und Schenkel. \hfill (Abitur 2023 GK)
\teil Zeige: Das Dreieck $A(0 | 3 | 1)$, $B(4 | 1 | 5)$, $C(4 | 6 | 3)$ ist gleichschenklig. Nenne Basis und Schenkel. \hfill (Abitur 2023 GK)
\teil Prüfe, ob das Dreieck $A(4 | 2 | 2)$, $B(2 | 4 | 2)$, $C(2 | 2 | 4)$ gleichseitig ist. \hfill (Abitur 2023 GK)
\teil Prüfe, ob das Dreieck $A(5 | 2 | 2)$, $B(2 | 5 | 2)$, $C(2 | 2 | 6)$ gleichseitig ist. \hfill (Abitur 2020 GK)
\teil Prüfe, ob das Dreieck $A(7 | 3 | 1)$, $B(3 | 7 | 1)$, $C(3 | 3 | 5)$ gleichseitig ist. \hfill (Abitur 2023 GK)
\teil $A(1 | 2 | 0)$, $B(5 | 6 | 0)$ und $C(3 | 4 | p)$ mit $p > 0$. Zeige: Das Dreieck $ABC$ ist für jedes $p$ gleichschenklig mit der Basis $AB$. \hfill (Abitur 2022 GK)
\teil $A(3 | 0 | 3)$, $B(3 | 6 | 1)$ und $C(p | 3 | 2)$. Zeige: Das Dreieck $ABC$ ist für jedes $p \ne 3$ gleichschenklig mit der Basis $AB$. \hfill (Abitur 2022 GK)
\teil $A(-1 | 2 | 2)$, $B(3 | 2 | -2)$ und $C(1 | p | 0)$. Zeige: Das Dreieck $ABC$ ist für jedes $p$ gleichschenklig. Welche Seite ist die Basis? \hfill (Abitur 2022 GK)
\teil Ein Quader hat die Ecken $A(9 | 0 | 0)$, $B(9 | 9 | 0)$ und $C(0 | 9 | 0)$ in der $x_1x_2$-Ebene. Begründe: Das Dreieck $ABC$ ist rechtwinklig und gleichschenklig. Flächeninhalt? \hfill (Abitur 2018 GK)
\end{teile}
\end{aufgabe}

\begin{aufgabe}{Dreiecke nachweisen – Prüfungsaufgaben – weiter}
\begin{teile}
\teil Grundfläche $A(6 | 1 | 0)$, $B(1 | 4 | 0)$, $C(1 | -2 | 0)$; Deckkante $E(1 | 3 | 2)$, $F(1 | -1 | 2)$. Zeige: $ABC$ ist gleichschenklig. Begründe: $EF$ ist parallel zur Grundfläche. \hfill (Abitur 2023 GK)
\teil Grundfläche $A(6 | 4 | 3)$, $B(2 | 1 | 3)$, $C(2 | 7 | 3)$; Deckkante $G(1 | 3 | 5)$, $H(1 | 5 | 5)$. Zeige: $ABC$ ist gleichschenklig. Begründe: $GH$ ist parallel zur Grundfläche. \hfill (Abitur 2023 GK)
\teil Grundfläche $A(5 | 6 | 0)$, $B(3 | 1 | 0)$, $C(7 | 1 | 0)$; Deckkante $P(4 | 2 | 3)$, $Q(6 | 2 | 3)$. Zeige: $ABC$ ist gleichschenklig. Begründe: $PQ$ ist parallel zur Grundfläche. \hfill (Abitur 2023 GK)
\teil $D(8 | 0 | 2)$, $E(0 | 6 | 2)$, $F(0 | 0 | 2)$ und $M(4 | 3 | 2)$. Begründe: $D$, $E$ und $F$ liegen auf einem Kreis mit dem Mittelpunkt $M$. \hfill (Abitur 2024 GK)
\teil Begründe: $D(2 | 2 | 5)$, $E(6 | 6 | 1)$ und $F(6 | 2 | 5)$ liegen auf einem Kreis, dessen Mittelpunkt die Mitte von $DE$ ist. \hfill (Abitur 2024 GK)
\teil Begründe: $P(-1 | 3 | 0)$, $Q(5 | 3 | 0)$ und $R(2 | 3 | 3)$ liegen auf einem Kreis, dessen Mittelpunkt die Mitte von $PQ$ ist. \hfill (Abitur 2024 GK)
\teil In der $x_1x_2$-Ebene liegen $A(0 | 0)$, $B(6 | 0)$ und $C(3 | 9)$ auf einem Kreis. Koordinaten seines Mittelpunkts $M$? \hfill (Abitur 2019 GK) \\ M( \leerfeld | \leerfeld )
\teil In der $x_1x_2$-Ebene liegen $A(-2 | 0)$, $B(6 | 0)$ und $C(2 | 8)$ auf einem Kreis. Koordinaten seines Mittelpunkts $M$? \hfill (Abitur 2019 GK) \\ M( \leerfeld | \leerfeld )
\teil In der $x_1x_2$-Ebene liegen $A(1 | 1)$, $B(7 | 1)$ und $C(4 | 10)$ auf einem Kreis. Koordinaten seines Mittelpunkts $M$? \hfill (Abitur 2019 GK) \\ M( \leerfeld | \leerfeld )
\end{teile}
\end{aufgabe}

\begin{aufgabe}{Dreiecke nachweisen – Prüfungsaufgaben – weiter}
\begin{teile}
\teil Dreieck $A(2 | 1 | 2)$, $B(4 | 5 | 4)$, $C(5 | 3 | 1)$ mit $|CA| = |CB|$. Gesucht ist ein Punkt $D \ne C$ mit $|DA| = |DB|$ und $|DA| = |CA|$. Koordinaten von $D$? \hfill (Abitur 2023 GK) \\ D( \leerfeld | \leerfeld | \leerfeld )
\teil Dreieck $A(0 | 2 | 3)$, $B(4 | 2 | -1)$, $C(4 | 3 | 3)$ mit $|CA| = |CB|$. Gesucht ist ein Punkt $D \ne C$ mit $|DA| = |DB|$ und $|DA| = |CA|$. Koordinaten von $D$? \hfill (Abitur 2023 GK) \\ D( \leerfeld | \leerfeld | \leerfeld )
\teil $A(3 | 1 | 2)$ und $M(1 | 3 | 1)$. Gib zwei Punkte $B$ und $C$ an, sodass $A$, $B$ und $C$ denselben Abstand zu $M$ haben und ein gleichschenkliges Dreieck bilden. \hfill (Abitur 2024 GK)
\teil $A(2 | 0 | 5)$ und $M(4 | 2 | 4)$. Gib zwei Punkte $B$ und $C$ an, sodass $A$, $B$ und $C$ denselben Abstand zu $M$ haben und ein gleichschenkliges Dreieck bilden. \hfill (Abitur 2024 GK)
\end{teile}
\end{aufgabe}

% Anlauf (ohne Prüfkennung)
\begin{aufgabe}{Vierecke nachweisen – Anlauf}
\begin{teile}
\teil Nachzuweisen: $ABCD$ ist ein Rechteck. Was gehört zum Steckbrief? \\ \kreuz{vier gleich lange Seiten} \\ \kreuz{Parallelogramm und ein rechter Winkel} \\ \kreuz{ein Paar paralleler Seiten}
\end{teile}
\end{aufgabe}

\begin{aufgabe}{Vierecke nachweisen – Prüfungsaufgaben}
\begin{teile}
\teil Zeige: $A(2 | 1 | 0)$, $B(4 | 3 | 1)$, $C(2 | 4 | 3)$, $D(0 | 2 | 2)$ bilden ein Rechteck $ABCD$. \hfill (Abitur 2023 GK)
\teil $A(1 | 5 | p)$, $B(4 | 5 | p)$, $C(4 | 1 | 5)$ und $D(1 | 1 | 5)$ mit einer natürlichen Zahl $p$. Zeige: $ABCD$ ist für jedes $p$ ein Rechteck. Länge von $BC$ in Abhängigkeit von $p$? \hfill (Abitur 2025 GK)
\teil Zeige: $A(0 | 1 | 3)$, $B(2 | 3 | 4)$, $C(3 | 5 | 6)$, $D(1 | 3 | 5)$ bilden eine Raute $ABCD$. \hfill (Abitur 2022 GK)
\teil Eine Dachfläche hat die Ecken $A(0 | 1 | 1)$, $B(4 | 1 | 4)$, $C(4 | 6 | 4)$, $D(0 | 6 | 1)$. Zeige: $ABCD$ ist eine Raute. \hfill (Abitur 2022 GK)
\teil Zeige: $A(1 | 2 | 2)$, $B(2 | 6 | 10)$, $C(6 | 7 | 18)$, $D(5 | 3 | 10)$ bilden eine Raute $ABCD$. \hfill (Abitur 2022 GK)
\teil Zeige: $A(2 | 2 | 2)$, $B(5 | 3 | 2)$, $C(6 | 5 | 4)$, $D(3 | 4 | 4)$ bilden ein Parallelogramm, aber kein Rechteck. \hfill (Abitur 2022 GK)
\teil Zeige: $A(1 | 3 | 2)$, $B(3 | 4 | 4)$, $C(5 | 2 | 5)$, $D(3 | 1 | 3)$ bilden eine Raute, aber kein Quadrat. \hfill (Abitur 2020 GK)
\teil Zeige: $A(1 | 0 | 1)$, $B(1 | 3 | 5)$, $C(4 | 7 | 5)$, $D(4 | 4 | 1)$ bilden eine Raute, aber kein Quadrat. \hfill (Abitur 2020 GK)
\teil Eine Kletterwand hat die Ecken $A(4 | 1 | 3)$, $B(1 | 4 | 3)$, $F(1 | 7 | 0)$ und $E(7 | 1 | 0)$ (1 LE = 1 m). Zeige: $ABFE$ ist ein Trapez, in dem zwei gegenüberliegende Seiten gleich lang sind. \hfill (Abitur 2018 GK)
\teil Ein Würfel mit Kante $8$ hat eine Ecke im Ursprung und liegt im ersten Oktanten. $P(8 | 0 | 2)$, $Q(3 | 8 | 0)$, $R(0 | 8 | 3)$ und $S(2 | 0 | 8)$ liegen auf seinen Kanten. Begründe: $PQRS$ ist ein Trapez mit zwei gleich langen Seiten, und $PS$ ist doppelt so lang wie $QR$. \hfill (Abitur 2019 GK)
\end{teile}
\end{aufgabe}

\begin{aufgabe}{Vierecke nachweisen – Prüfungsaufgaben – weiter}
\begin{teile}
\teil Eine Terrasse ist das Viereck $A(1 | -3 | 0)$, $B(7 | -1 | 0)$, $C(7 | 1 | 0)$, $D(1 | 3 | 0)$ (1 LE = 1 m). Zeige: $ABCD$ ist ein Trapez. Flächeninhalt? \hfill (Abitur 2026 GK) \\ \leerfeld[m²]
\teil Ein Beet ist das Viereck $E(-2 | 1 | 0)$, $F(6 | 1 | 0)$, $G(3 | 5 | 0)$, $H(1 | 5 | 0)$ (1 LE = 1 m). Zeige: $EFGH$ ist ein Trapez. Flächeninhalt? \hfill (Abitur 2026 GK) \\ \leerfeld[m²]
\teil Eine Glasplatte ist das Viereck $K(0 | -3 | 1)$, $L(4 | -1 | 1)$, $M(4 | 1 | 1)$, $N(0 | 3 | 1)$ (1 LE = 1 dm). Zeige: $KLMN$ ist ein Trapez. Flächeninhalt? \hfill (Abitur 2026 GK) \\ \leerfeld dm²
\teil Zeige: Das Viereck $K(0 | 0 | 2)$, $L(2 | 3 | 2)$, $M(6 | 0 | 2)$, $N(2 | -3 | 2)$ ist ein Drachenviereck. \hfill (Abitur 2023 GK)
\teil Zeige: Das Viereck $K(1 | 2 | 0)$, $L(1 | 5 | 2)$, $M(1 | 2 | 8)$, $N(1 | -1 | 2)$ ist ein Drachenviereck. \hfill (Abitur 2023 GK)
\teil Zeige: Das Viereck $P(6 | 1 | 1)$, $Q(6 | 5 | 1)$, $R(6 | 4 | 4)$, $S(6 | 1 | 5)$ ist ein Drachenviereck. \hfill (Abitur 2023 GK)
\teil Eine ebene Rasenfläche hat die Ecken $A(2 | 1 | 0)$, $B(9 | 1 | 0)$, $C(8 | 4 | 0{,}5)$, $D(8 | 7 | 1)$ und $E(2 | 7 | 1)$ (1 LE = 1 m). Zeige: $AE$ ist parallel zu $CD$, und $CD$ und $DE$ bilden einen rechten Winkel. \hfill (Abitur 2021 GK)
\teil Zeige: Im Viereck $A(1 | -1 | 4)$, $B(3 | 1 | 5)$, $C(1 | 2 | 2{,}5)$, $D(0 | 1 | 2)$ ist $AB$ parallel zu $DC$, und bei $A$ liegt ein rechter Winkel. \hfill (Abitur 2021 GK)
\teil Zeige: Im Viereck $A(0 | 2 | 1)$, $B(6 | 2 | 9)$, $C(7 | 3 | 2)$, $D(4 | 3 | -2)$ ist $AB$ parallel zu $DC$, und bei $A$ liegt ein rechter Winkel. \hfill (Abitur 2021 GK)
\teil $P(2 | 1 | 0)$ und $Q(2 | 1 | 5)$ sind benachbarte Ecken eines Quadrats in der Ebene $3x_1 + 4x_2 = 10$. Gib eine weitere Ecke an. \hfill (Abitur 2022 GK)
\teil $P(3 | 1 | 4)$ und $Q(3 | 6 | 4)$ sind benachbarte Ecken eines Quadrats in der Ebene $4x_1 - 3x_3 = 0$. Gib eine weitere Ecke an. \hfill (Abitur 2022 GK)
\end{teile}
\end{aufgabe}

\begin{aufgabe}{Vierecke nachweisen – Prüfungsaufgaben – weiter}
\begin{teile}
\teil Das Viereck $ABCD$ hat die Ecken $A(0 | 0 | 0)$, $B(3 | 1 | 0)$, $C(0 | 7 | 0)$ und $D(-2 | 3 | 0)$. $B$ wird parallel zur $x_2$-Achse nach $B'$ verschoben, sodass $AB'CD$ bei $B'$ einen rechten Innenwinkel hat; $b$ ist die gesuchte Koordinate. Erläutere den Ansatz $(3 | b | 0) \circ (3 | b - 7 | 0) = 0$. \hfill (Abitur 2025 GK)
\teil Das Viereck $ABCD$ hat die Ecken $A(0 | 0 | 0)$, $B(2 | 4 | 0)$, $C(9 | 0 | 0)$ und $D(4 | -3 | 0)$. $B$ wird parallel zur $x_1$-Achse nach $B'$ verschoben, sodass $AB'CD$ bei $B'$ einen rechten Innenwinkel hat; $b$ ist die gesuchte Koordinate. Erläutere den Ansatz $(b | 4 | 0) \circ (b - 9 | 4 | 0) = 0$. \hfill (Abitur 2025 GK)
\teil Die Raute $A(5 | 3 | 1)$, $B(2 | 3 | 5)$, $C(-1 | 3 | 1)$, $D(2 | 3 | -3)$ hat einen Inkreis, der alle vier Seiten berührt. Begründe: $P(3{,}92 | 3 | 2{,}44)$ ist einer der Berührpunkte. \hfill (Abitur 2020 GK)
\teil Die Raute $A(1 | 1 | 6)$, $B(4 | 1 | 2)$, $C(1 | 1 | -2)$, $D(-2 | 1 | 2)$ hat einen Inkreis, der alle vier Seiten berührt. Begründe: $P(2{,}92 | 1 | 3{,}44)$ ist einer der Berührpunkte. \hfill (Abitur 2020 GK)
\teil Eine Wand hat die Form des Parallelogramms $A(2 | 1 | 0)$, $B(9 | 1 | 0)$, $C(10 | 1 | 4)$, $D(3 | 1 | 4)$ (1 LE = 1 m). Ein Farbstreifen von $A$ nach $P(6 | 1 | 4)$ auf der Seite $DC$ teilt sie. Welcher Teil ist größer? Begründe ohne Flächenberechnung. \hfill (Abitur 2022 GK)
\teil Eine rechteckige Wand hat die Ecken $A(1 | 3 | 0)$, $B(1 | 9 | 0)$, $C(1 | 9 | 5)$, $D(1 | 3 | 5)$ (1 LE = 1 m). Ein Streifen von $B$ nach $Q(1 | 3 | 2)$ auf der Seite $AD$ teilt sie. Welcher Teil ist größer? Begründe ohne Flächenberechnung. \hfill (Abitur 2022 GK)
\end{teile}
\end{aufgabe}

% Anlauf (ohne Prüfkennung)
\begin{aufgabe}{Körper und Drehungen – Anlauf}
\begin{teile}
\teil Ein Quader steht auf der $x_1x_2$-Ebene; seine Deckenecke $G$ hat $x_3 = 4$. Was sagt das über die Deckfläche? \\ \kreuz{Sie liegt in der $x_1x_2$-Ebene.} \\ \kreuz{Sie liegt parallel zur $x_1x_2$-Ebene in der Höhe $4$.} \\ \kreuz{Sie steht senkrecht auf der $x_1x_2$-Ebene.}
\end{teile}
\end{aufgabe}

\begin{aufgabe}{Körper und Drehungen – Prüfungsaufgaben}
\begin{teile}
\teil Gerades Prisma $ABCDEF$ ($D$ über $A$, $E$ über $B$, $F$ über $C$) mit $A(2 | 1 | 0)$, $B(7 | 1 | 0)$, $C(2 | 4 | 0)$, $D(2 | 1 | 3)$: $F$? \hfill (Abitur 2026 GK) \\ F( \leerfeld | \leerfeld | \leerfeld )
\teil Quader $ABCDEFGH$ (Boden $ABCD$, $E$ über $A$, $F$ über $B$, $G$ über $C$, $H$ über $D$) mit $A(2 | 1 | 0)$, $B(6 | 1 | 0)$, $E(2 | 1 | 5)$, $H(2 | 4 | 5)$: $G$? \hfill (Abitur 2024 GK) \\ G( \leerfeld | \leerfeld | \leerfeld )
\teil Quader $ABCDEFGH$ (Boden $ABCD$, $E$ über $A$, $F$ über $B$, $G$ über $C$, $H$ über $D$) mit $A(0 | 2 | 1)$, $B(4 | 2 | 1)$, $E(0 | 2 | 4)$, $H(0 | 7 | 4)$: $G$? \hfill (Abitur 2024 GK) \\ G( \leerfeld | \leerfeld | \leerfeld )
\teil Gerades Prisma $ABCDEF$ ($D$ über $A$, $E$ über $B$, $F$ über $C$) mit $A(4 | 1 | 2)$, $B(4 | 6 | 2)$, $C(8 | 1 | 2)$, $D(4 | 1 | 7)$: $F$? \hfill (Abitur 2026 GK) \\ F( \leerfeld | \leerfeld | \leerfeld )
\teil Eine Pyramide hat als Grundfläche ein rechtwinkliges Dreieck mit den Katheten $6$ cm und $8$ cm; die Spitze liegt $5$ cm senkrecht über der Ecke mit dem rechten Winkel. Wähle geeignete Koordinaten der vier Ecken (1 LE = 1 cm). \hfill (Abitur 2020 GK)
\teil Eine Halle ist ein Quader mit quadratischem Grundriss (Seite $15$ m) und $7$ m Höhe. $A$ liegt im Ursprung, $B$ auf der $x_1$-Achse, $D$ auf der $x_2$-Achse, die Deckenecken $E$, $F$, $G$, $H$ über $A$, $B$, $C$, $D$ (1 LE = 1 m). Koordinaten von $G$ und $H$? \hfill (Abitur 2020 GK)
\teil Ein Würfel mit Kante $2$ hat eine Ecke im Ursprung, alle Ecken haben Koordinaten $\ge 0$. Er wird um $(-1 | -1 | 1)$ verschoben. Welche Ecke hat danach $x_1 < 0$, $x_2 > 0$ und $x_3 > 2$? \hfill (Abitur 2024 GK) \\ ( \leerfeld | \leerfeld | \leerfeld )
\teil Ein Quader mit den Kanten $4$, $2$ und $3$ entlang der $x_1$-, $x_2$- und $x_3$-Achse hat eine Ecke im Ursprung, alle Ecken haben Koordinaten $\ge 0$. Er wird um $(-2 | 1 | -1)$ verschoben. Welche Ecke hat danach $x_1 > 0$, $x_2 < 2$ und $x_3 < 0$? \hfill (Abitur 2024 GK) \\ ( \leerfeld | \leerfeld | \leerfeld )
\teil Ein Quader mit den Kanten $3$, $5$ und $2$ entlang der $x_1$-, $x_2$- und $x_3$-Achse hat eine Ecke im Ursprung, alle Ecken haben Koordinaten $\ge 0$. Er wird um $(-2 | -3 | -1)$ verschoben. Welche Ecke hat danach $x_1 > 0$, $x_2 < 0$ und $x_3 > 0$? \hfill (Abitur 2024 GK) \\ ( \leerfeld | \leerfeld | \leerfeld )
\teil Eine Pyramide hat die Grundfläche $A(4 | -1 | 2)$, $B(1 | 3 | 2)$, $C(-1 | 0 | 2)$ und die Spitze $S(2 | 1 | 7)$. Begründe: Die Grundfläche ist parallel zur $x_1x_2$-Ebene. Höhe der Pyramide? \hfill (Abitur 2018 GK) \\ h = \leerfeld
\teil Eine Pyramide hat die Grundfläche $A(2 | 2 | -1)$, $B(5 | 0 | -1)$, $C(3 | 5 | -1)$ und die Spitze $S(3 | 2 | 3)$. Begründe: Die Grundfläche ist parallel zur $x_1x_2$-Ebene. Höhe der Pyramide? \hfill (Abitur 2018 GK) \\ h = \leerfeld
\teil Eine Pyramide hat die Grundfläche $A(1 | 0 | 3)$, $B(4 | 2 | 3)$, $C(0 | 5 | 3)$ und die Spitze $S(2 | 2 | 9)$. Begründe: Die Grundfläche ist parallel zur $x_1x_2$-Ebene. Höhe der Pyramide? \hfill (Abitur 2018 GK) \\ h = \leerfeld
\end{teile}
\end{aufgabe}

\begin{aufgabe}{Körper und Drehungen – Prüfungsaufgaben – weiter}
\begin{teile}
\teil Ein Dachgeschoss hat einen rechteckigen Boden von $10$ m mal $8$ m. Die Dachflächen steigen von den beiden Traufwänden (Höhe $0{,}5$ m) gleichmäßig zum First über der Mitte (Höhe $5{,}5$ m). Flächen mit weniger als $2$ m Raumhöhe zählen nicht zur Wohnfläche. Welcher Anteil des Bodens zählt nicht? \hfill (Abitur 2019 GK) \\ \leerfeld \%
\teil Ein Dachgeschoss hat einen rechteckigen Boden von $14$ m mal $10$ m. Die Dachflächen steigen von den beiden Traufwänden (Höhe $1$ m) gleichmäßig zum First über der Mitte (Höhe $5$ m). Flächen mit weniger als $1{,}8$ m Raumhöhe zählen nicht zur Wohnfläche. Welcher Anteil des Bodens zählt nicht? \hfill (Abitur 2019 GK) \\ \leerfeld \%
\teil Ein Dachgeschoss hat einen rechteckigen Boden von $15$ m mal $10$ m. Die Dachflächen steigen von den beiden Traufwänden (Höhe $0{,}4$ m) gleichmäßig zum First über der Mitte (Höhe $2{,}9$ m). Flächen mit weniger als $1$ m Raumhöhe zählen nicht zur Wohnfläche. Welcher Anteil des Bodens zählt nicht? \hfill (Abitur 2019 GK) \\ \leerfeld \%
\teil Das Dreieck $A(5 | 3 | 0)$, $B(3 | 5 | 0)$, $C(3 | 3 | 4)$ wird um die Kante $AB$ gedreht, bis $C$ in der $x_1x_2$-Ebene liegt, mit $x_1 < 4$. Koordinaten des gedrehten Punktes $C'$? \hfill (Abitur 2023 GK) \\ C'( \leerfeld | \leerfeld | \leerfeld )
\teil Das Dreieck $A(2 | 1 | 0)$, $B(2 | 9 | 0)$, $C(5 | 5 | 4)$ wird um die Kante $AB$ gedreht, bis $C$ in der $x_1x_2$-Ebene liegt, mit $x_1 > 2$. Koordinaten des gedrehten Punktes $C'$? \hfill (Abitur 2023 GK) \\ C'( \leerfeld | \leerfeld | \leerfeld )
\teil Eine Pyramide hat den quadratischen Boden $ABCD$ mit der Seite $4$ in der $x_1x_2$-Ebene, $A$ im Ursprung, $B$ auf der $x_1$-Achse, $D$ auf der $x_2$-Achse; die Spitze $S$ liegt $6$ senkrecht über $A$. Ein Würfel hat drei Seitenflächen in den Koordinatenebenen und eine Ecke auf der Kante $CS$. Kantenlänge des Würfels? Begründe, dass kein Punkt des Würfels außerhalb der Pyramide liegt. \hfill (Abitur 2024 GK)
\teil Eine Pyramide hat den quadratischen Boden $ABCD$ mit der Seite $3$ in der $x_1x_2$-Ebene, $A$ im Ursprung, $B$ auf der $x_1$-Achse, $D$ auf der $x_2$-Achse; die Spitze $S$ liegt $6$ senkrecht über $A$. Ein Würfel hat drei Seitenflächen in den Koordinatenebenen und eine Ecke auf der Kante $CS$. Kantenlänge des Würfels? Begründe, dass kein Punkt des Würfels außerhalb der Pyramide liegt. \hfill (Abitur 2024 GK)
\teil Ein Körper hat den rechteckigen Boden $A(2 | 1 | 0)$, $B(6 | 1 | 0)$, $C(6 | 7 | 0)$, $D(2 | 7 | 0)$ und die Firstkante $P(4 | 2 | 3)$, $Q(4 | 5 | 3)$; die Dreiecke $ABP$ und $CDQ$ und die Vierecke $BCQP$ und $DAPQ$ begrenzen ihn. Die Ebenen $x_2 = s$ mit $1 < s < 7$ schneiden ihn. Wie viele Ecken hat die Schnittfläche? Für welche $s$ haben die Schnittflächen dieselbe Form und denselben Flächeninhalt? \hfill (Abitur 2024 GK)
\teil Ein Körper hat den rechteckigen Boden $A(0 | 1 | 0)$, $B(6 | 1 | 0)$, $C(6 | 9 | 0)$, $D(0 | 9 | 0)$ und die Firstkante $P(3 | 3 | 4)$, $Q(3 | 6 | 4)$; die Dreiecke $ABP$ und $CDQ$ und die Vierecke $BCQP$ und $DAPQ$ begrenzen ihn. Die Ebenen $x_2 = s$ mit $1 < s < 9$ schneiden ihn. Wie viele Ecken hat die Schnittfläche? Für welche $s$ haben die Schnittflächen dieselbe Form und denselben Flächeninhalt? \hfill (Abitur 2024 GK)
\teil Ein Würfel $ABCDEFGH$ mit Kante $4$ hat die Ecke $A$ im Ursprung, $B$ auf der $x_1$-Achse, $D$ auf der $x_2$-Achse; $E$ liegt über $A$, $G$ über $C$. Die Ebene durch $E$, $G$ und $S_t(t | 0 | 0)$ mit $0 \le t \le 4$ schneidet ihn. Eckenzahl der Schnittfläche in Abhängigkeit von $t$? Für welche $t$ hat sie genau zwei Symmetrieachsen? \hfill (Abitur 2024 GK)
\teil Ein Quader $ABCDEFGH$ mit quadratischem Boden (Seite $6$) und Höhe $3$ hat die Ecke $A$ im Ursprung, $B$ auf der $x_1$-Achse, $D$ auf der $x_2$-Achse; $E$ liegt über $A$, $G$ über $C$. Die Ebene durch $E$, $G$ und $S_t(t | 0 | 0)$ mit $0 \le t \le 6$ schneidet ihn. Eckenzahl der Schnittfläche in Abhängigkeit von $t$? Für welche $t$ hat sie genau zwei Symmetrieachsen? \hfill (Abitur 2024 GK)
\end{teile}
\end{aufgabe}

\clearpage
% Vektoren und Rechenoperationen – bank/vektoren-und-rechenoperationen/ (zusammenbau v0.4, Heft)
\einheitenkopf[t2]{Vektoren und Rechenoperationen}
\setcounter{aufgabe}{18}

% Anlauf (ohne Prüfkennung)
\begin{aufgabe}{Vektorbegriff und Betrag – Anlauf}
\begin{teile}
\teil Die Länge der Strecke von $A$ nach $B$: Zahl oder Vektor? \\ \kreuz{Zahl} \\ \kreuz{Vektor}
\teil $A(1 | 2 | 0)$, $B(4 | 3 | 2)$: $\overrightarrow{AB}$, $\overrightarrow{BA}$ und $2 \cdot \overrightarrow{AB}$?
\teil $\vec v = (1 | a | 2)$, $a < 0$, $|\vec v| = 3$: $a$? a = \leerfeld
\end{teile}
\end{aufgabe}

\begin{aufgabe}{Vektorbegriff und Betrag – Prüfungsaufgaben}
\begin{teile}
\teil $A(1 | 2 | a)$ und $B(6 | 14 | 3)$ haben den Abstand 13: $a$? \hfill (Abitur 2021 GK) \\ a = \leerfeld
\teil Ein Verleih verteilt 600 Leihräder auf drei Stationen; jede Verteilung wird als Vektor (Räder an Station 1 | 2 | 3) notiert. Können zwei solche Vektoren verschiedene Beträge haben? Begründe mit einem Beispiel. \hfill (Abitur 2024 GK)
\teil Bei einer Abstimmung mit drei Listen werden 90 Stimmen verteilt; jedes Ergebnis wird als Vektor (Stimmen für Liste 1 | 2 | 3) notiert. Haben alle Ergebnisvektoren denselben Betrag? Begründe mit einem Beispiel. \hfill (Abitur 2024 GK)
\end{teile}
\end{aufgabe}

% Anlauf (ohne Prüfkennung)
\begin{aufgabe}{Vektorterme am Körper – Anlauf}
\begin{teile}
\teil Wohin führt $\overrightarrow{AP} = \overrightarrow{AB} + \overrightarrow{BF}$? \\ \kreuz{Ecke} \\ \kreuz{Kantenmitte} \\ \kreuz{Flächenmitte}

\begin{ksys3}[x1max=7] \rquader{0,0,0}{6}{4}{3} \rpunkt*{0}{0}{0}{A} \rpunkt*{6}{0}{0}{B} \rpunkt*{6}{4}{0}{C} \rpunkt*{0}{4}{0}{D} \rpunkt*{0}{0}{3}{E} \rpunkt*{6}{0}{3}{F} \rpunkt*{6}{4}{3}{G} \rpunkt*{0}{4}{3}{H} \end{ksys3}
\teil Trage $P$ mit $\overrightarrow{AP} = \overrightarrow{AD} + \frac{1}{2} \cdot \overrightarrow{DC}$ ein.

\begin{ksys3}[x1max=7] \rquader{0,0,0}{6}{4}{3} \rpunkt*{0}{0}{0}{A} \rpunkt*{6}{0}{0}{B} \rpunkt*{6}{4}{0}{C} \rpunkt*{0}{4}{0}{D} \rpunkt*{0}{0}{3}{E} \rpunkt*{6}{0}{3}{F} \rpunkt*{6}{4}{3}{G} \rpunkt*{0}{4}{3}{H} \end{ksys3}
\teil $M_1$ Mitte von $AB$, $M_2$ Mitte von $GH$: Zeichne $\overrightarrow{M_1M_2}$ ein. $r$, $s$, $t$ mit $\overrightarrow{M_1M_2} = r\vec u + s\vec v + t\vec w$?

\begin{ksys3}[x1max=7] \rquader{0,0,0}{6}{4}{3} \rpunkt*{0}{0}{0}{A} \rpunkt*{6}{0}{0}{B} \rpunkt*{6}{4}{0}{C} \rpunkt*{0}{4}{0}{D} \rpunkt*{0}{0}{3}{E} \rpunkt*{6}{0}{3}{F} \rpunkt*{6}{4}{3}{G} \rpunkt*{0}{4}{3}{H} \rvektorab[below]{0,0,0}{6,0,0}{\vec u} \rvektorab[left]{0,0,0}{0,4,0}{\vec v} \rvektorab[left]{0,0,0}{0,0,3}{\vec w} \end{ksys3}
r = \leerfeld , s = \leerfeld , t = \leerfeld
\end{teile}
\end{aufgabe}

\begin{aufgabe}{Vektorterme am Körper – Prüfungsaufgaben}
\begin{teile}
\teil Trage $P$ mit $\overrightarrow{AP} = \frac{1}{2} \cdot \overrightarrow{AC} + \overrightarrow{AE}$ ein. \hfill (Abitur 2021 GK)

\begin{ksys3}[x1max=7] \rquader{0,0,0}{6}{4}{3} \rpunkt*{0}{0}{0}{A} \rpunkt*{6}{0}{0}{B} \rpunkt*{6}{4}{0}{C} \rpunkt*{0}{4}{0}{D} \rpunkt*{0}{0}{3}{E} \rpunkt*{6}{0}{3}{F} \rpunkt*{6}{4}{3}{G} \rpunkt*{0}{4}{3}{H} \end{ksys3}
\teil Trage $P$ mit $\overrightarrow{AP} = \overrightarrow{AB} + \frac{1}{2} \cdot \overrightarrow{BC} + \overrightarrow{CG}$ ein. \hfill (Abitur 2026 GK)

\begin{ksys3}[x1max=7] \rquader{0,0,0}{6}{4}{3} \rpunkt*{0}{0}{0}{A} \rpunkt*{6}{0}{0}{B} \rpunkt*{6}{4}{0}{C} \rpunkt*{0}{4}{0}{D} \rpunkt*{0}{0}{3}{E} \rpunkt*{6}{0}{3}{F} \rpunkt*{6}{4}{3}{G} \rpunkt*{0}{4}{3}{H} \end{ksys3}
\teil Beschreibe die Lage von $P$ mit $\overrightarrow{AP} = \frac{1}{2} \cdot (\overrightarrow{AG} - \overrightarrow{AB})$. \hfill (Abitur 2021 GK)

\begin{ksys3}[x1max=7] \rquader{0,0,0}{6}{4}{3} \rpunkt*{0}{0}{0}{A} \rpunkt*{6}{0}{0}{B} \rpunkt*{6}{4}{0}{C} \rpunkt*{0}{4}{0}{D} \rpunkt*{0}{0}{3}{E} \rpunkt*{6}{0}{3}{F} \rpunkt*{6}{4}{3}{G} \rpunkt*{0}{4}{3}{H} \end{ksys3}
\teil $A(2 | 1 | 0)$, $B(4 | 2 | 2)$, $C(3 | 4 | 2)$; das Dreieck $ABC$ hat bei $B$ einen rechten Winkel. $D$ so, dass $ABCD$ ein Rechteck ist? \hfill (Abitur 2017 GK) \\ D( \leerfeld | \leerfeld | \leerfeld )
\teil Quader mit $A(2 | 0 | 1)$, $G(6 | 4 | 5)$, $H(2 | 4 | 5)$; er wird so verschoben, dass der Schnittpunkt seiner Raumdiagonalen im Ursprung liegt: $H'$? \hfill (Abitur 2024 GK) \\ H'( \leerfeld | \leerfeld | \leerfeld )
\end{teile}
\end{aufgabe}

\begin{aufgabe}{Vektorterme am Körper – Prüfungsaufgaben – weiter}
\begin{teile}
\teil Ein Prisma steht mit der Kante $AB$ von $A(3 | 0 | 0)$ nach $B(0 | 4 | 0)$ auf der $x_1x_2$-Ebene; $F(0 | 4 | 2)$ liegt senkrecht über $B$. Es wird um $AB$ gedreht, bis $F$ in der $x_1x_2$-Ebene liegt und eine positive $x_1$-Koordinate hat: $F'$. Term für $\overrightarrow{OF'}$? \hfill (Abitur 2026 GK)
\teil Eine Klappe steht senkrecht mit der Kante $AB$ von $A(4 | 1 | 0)$ nach $B(0 | 4 | 0)$ auf der $x_1x_2$-Ebene; $F(0 | 4 | 4)$ liegt senkrecht über $B$. Sie wird um $AB$ gedreht, bis $F$ in der $x_1x_2$-Ebene liegt und eine positive $x_1$-Koordinate hat: $F'$. Term für $\overrightarrow{OF'}$? \hfill (Abitur 2026 GK)
\end{teile}
\end{aufgabe}

% Anlauf (ohne Prüfkennung)
\begin{aufgabe}{Skalarprodukt im Sachzusammenhang – Anlauf}
\begin{teile}
\teil Die Liste der gekauften Mengen je Sorte (Äpfel, Birnen, Pflaumen): Zahl oder Vektor? \\ \kreuz{Zahl} \\ \kreuz{Vektor}
\teil Ein Café verkauft an einem Tag 12 Espresso, 30 Cappuccino und 18 Tee. Absatzvektor (Espresso | Cappuccino | Tee)? Was zählt die zweite Komponente? ( \leerfeld | \leerfeld | \leerfeld )
\teil Ein Imbiss kauft $\vec m = (3 | 2 | 5)$ Kisten, die Preise sind $\vec p = (14 | 18 | 6)$ Euro je Kiste. $\vec m \circ \vec p$ und Bedeutung?
\end{teile}
\end{aufgabe}

\begin{aufgabe}{Skalarprodukt im Sachzusammenhang – Prüfungsaufgaben}
\begin{teile}
\teil Ein Imbiss kauft $\vec m = (k | s | w)$ Kisten Cola, Saft und Wasser; $\vec p = (14{,}50 | 18{,}00 | 6{,}40)$ sind die Preise in Euro je Kiste. Was gibt $\vec m \circ \vec p$ an? \hfill (Abitur 2019 GK)
\teil Ein Maler mischt Weiß, Gelb und Rot im Verhältnis 5 : 2 : 1; $\vec m = (w | g | r)$ in Litern, $\vec p = (4 | 6 | 8)$ Euro je Liter, $\vec m \circ \vec p = 168$. Wie viel Liter jeder Farbe? \hfill (Abitur 2019 GK)
\teil Ein Müsli besteht aus Haferflocken, Nüssen und Rosinen im Verhältnis 6 : 1 : 2; $\vec m = (h | n | c)$ in kg, $\vec p = (3 | 12 | 5)$ Euro je kg, $\vec m \circ \vec p = 280$. Wie viel kg jeder Zutat? \hfill (Abitur 2019 GK)
\end{teile}
\end{aufgabe}

\clearpage
% Geraden – bank/geraden/ (zusammenbau v0.4, Heft)
\einheitenkopf[t3]{Geraden}
\setcounter{aufgabe}{25}

% Anlauf (ohne Prüfkennung)
\begin{aufgabe}{Geradengleichung aufstellen und lesen – Anlauf}
\begin{teile}
\teil Gegeben ist $g\colon \vec x = (4 | -1 | 2) + t \cdot (3 | 0 | -2)$. Welcher Vektor gibt die Richtung von $g$ an? \kreuz{$(4 | -1 | 2)$} \\ \kreuz{$(3 | 0 | -2)$} \\ \kreuz{$(7 | -1 | 0)$}
\teil Stelle eine Gleichung der Geraden $g$ durch $A(2 | 1 | 3)$ und $B(4 | 4 | 2)$ auf.
\end{teile}
\end{aufgabe}

\begin{aufgabe}{Geradengleichung aufstellen und lesen – Prüfungsaufgaben}
\begin{teile}
\teil Ein Skilift führt geradlinig von der Talstation $T(2 | -4 | 1)$ zur Bergstation $B(-6 | 20 | 13)$; 1 LE entspricht 10 m. Gib eine Gleichung der Liftstrecke an und beschreibe, was sie im Sachzusammenhang darstellt. \hfill (Abitur 2023 GK)
\teil Eine Seilrutsche wird durch $\vec x = (0 | 5 | 12) + \lambda \cdot (15 | -10 | -10)$ mit $\lambda \in [0; 1]$ beschrieben; 1 LE entspricht 1 m, die $x_1x_2$-Ebene ist der Boden. Beschreibe die Bedeutung der Gleichung im Sachzusammenhang. \hfill (Abitur 2023 GK)
\teil Ein Förderband wird durch $\vec x = (3 | 1 | 0) + \mu \cdot (6 | 3 | 6)$ mit $0 \le \mu \le 1$ beschrieben; 1 LE entspricht 1 m, die $x_1x_2$-Ebene ist der Boden. Beschreibe die Bedeutung der Gleichung im Sachzusammenhang. \hfill (Abitur 2023 GK)
\teil Gegeben ist $g\colon \vec x = (3 | 1 | 2) + r \cdot (1 | 3 | -2)$, $r \in \mathbb{R}$. Gib eine Gleichung einer Geraden $p$ an, die echt parallel zu $g$ verläuft, und eine Gleichung einer Geraden $q$, die $g$ senkrecht schneidet. \hfill (Abitur 2019 GK)
\teil Gegeben ist $g\colon \vec x = (-2 | 4 | 1) + t \cdot (2 | -1 | 2)$, $t \in \mathbb{R}$. Gib eine Gleichung einer Geraden $p$ an, die echt parallel zu $g$ verläuft, und eine Gleichung einer Geraden $q$, die $g$ senkrecht schneidet. \hfill (Abitur 2019 GK)
\teil Gegeben ist $g\colon \vec x = (1 | -3 | 5) + r \cdot (3 | 1 | -1)$, $r \in \mathbb{R}$. Gib eine Gleichung einer Geraden $p$ an, die echt parallel zu $g$ verläuft, und eine Gleichung einer Geraden $q$, die $g$ senkrecht schneidet. \hfill (Abitur 2019 GK)
\teil Die Geraden $g$ und $h$ verlaufen parallel mit dem Richtungsvektor $(1 | 4 | 2)$; $P(2 | -1 | 3)$ liegt auf $g$, $Q(10 | 3 | -1)$ auf $h$. Der Punkt $T$ liegt auf der Strecke $PQ$ mit $|TQ| = 3 \cdot |PT|$. Gib eine Gleichung der Geraden an, die durch $T$ parallel zu $g$ und $h$ verläuft. \hfill (Abitur 2022 GK)
\teil Die Geraden $g$ und $h$ verlaufen parallel mit dem Richtungsvektor $(2 | -1 | 1)$; $P(-3 | 0 | 5)$ liegt auf $g$, $Q(3 | 6 | -1)$ auf $h$. Der Punkt $T$ liegt auf der Strecke $PQ$ mit $|TQ| = 2 \cdot |PT|$. Gib eine Gleichung der Geraden an, die durch $T$ parallel zu $g$ und $h$ verläuft. \hfill (Abitur 2022 GK)
\teil Die Geraden $g$ und $h$ verlaufen parallel mit dem Richtungsvektor $(1 | 1 | -2)$; $P(0 | 4 | 1)$ liegt auf $g$, $Q(8 | 0 | 5)$ auf $h$. Der Punkt $T$ liegt auf der Strecke $PQ$ mit $|PT| = 3 \cdot |TQ|$. Gib eine Gleichung der Geraden an, die durch $T$ parallel zu $g$ und $h$ verläuft. \hfill (Abitur 2022 GK)
\teil Der Punkt $P(3 | 2 | 2)$ liegt in der Ebene $E\colon \vec x = (2 | -1 | 1) + r \cdot (1 | 0 | 2) + t \cdot (0 | 3 | -1)$. Die Gerade $g$ liegt in $E$, geht durch $P$ und hat keinen Schnittpunkt mit der $x_1x_2$-Ebene. Gib eine Gleichung von $g$ an. \hfill (Abitur 2021 GK)
\teil Der Punkt $P(4 | 1 | 6)$ liegt in der Ebene $E\colon \vec x = (0 | 2 | 3) + r \cdot (2 | 1 | -1) + t \cdot (1 | -1 | 2)$. Die Gerade $g$ liegt in $E$, geht durch $P$ und hat keinen Schnittpunkt mit der $x_1x_2$-Ebene. Gib eine Gleichung von $g$ an. \hfill (Abitur 2021 GK)
\end{teile}
\end{aufgabe}

% Anlauf (ohne Prüfkennung)
\begin{aufgabe}{Punktprobe und Punkte auf der Geraden – Anlauf}
\begin{teile}
\teil Für $g\colon \vec x = (2 | 5 | 1) + t \cdot (0 | 3 | 4)$ und $P(2 | 11 | 9)$: Aus welcher Koordinate lässt sich $t$ bestimmen? \kreuz{nur aus der ersten} \\ \kreuz{aus der zweiten oder der dritten} \\ \kreuz{aus keiner}
\teil Prüfe, ob $P(7 | 5 | 8)$ auf $g\colon \vec x = (1 | 2 | -1) + r \cdot (2 | 1 | 3)$ liegt.
\end{teile}
\end{aufgabe}

\begin{aufgabe}{Punktprobe und Punkte auf der Geraden – Prüfungsaufgaben}
\begin{teile}
\teil Gegeben sind $g\colon \vec x = (5 | -1 | 2) + t \cdot (3 | 0 | -2)$, $t \in \mathbb{R}$, und $A(2 | 3 | 4)$. Begründe, dass $A$ nicht auf $g$ liegt. \hfill (Abitur 2023 GK)
\teil Gegeben sind $h\colon \vec x = (0 | 4 | 6) + r \cdot (2 | -3 | 0)$, $r \in \mathbb{R}$, und $B(4 | -2 | 5)$. Begründe, dass $B$ nicht auf $h$ liegt. \hfill (Abitur 2023 GK)
\teil Gegeben sind $k\colon \vec x = (-3 | 2 | 1) + t \cdot (0 | 1 | 4)$, $t \in \mathbb{R}$, und $C(1 | 5 | 13)$. Begründe, dass $C$ nicht auf $k$ liegt. \hfill (Abitur 2023 GK)
\teil Der Punkt $B(-2 | y_B | z_B)$ liegt auf $g\colon \vec x = (1 | 3 | 2) + r \cdot (-1 | 2 | -3)$. Bestimme $y_B$ und $z_B$. \hfill (Abitur 2019 GK) \\ $y_B$ = \leerfeld , $z_B$ = \leerfeld
\teil Der Punkt $P(x_P | 2 | z_P)$ liegt auf $g\colon \vec x = (0 | -1 | 4) + t \cdot (2 | 1 | -2)$. Bestimme $x_P$ und $z_P$. \hfill (Abitur 2019 GK) \\ $x_P$ = \leerfeld , $z_P$ = \leerfeld
\teil Der Punkt $Q(x_Q | y_Q | -3)$ liegt auf $h\colon \vec x = (3 | -2 | 0) + t \cdot (-2 | 4 | 1)$. Bestimme $x_Q$ und $y_Q$. \hfill (Abitur 2019 GK) \\ $x_Q$ = \leerfeld , $y_Q$ = \leerfeld
\teil Die Gerade $g$ verläuft durch $A(3 | 1 | 2)$ mit dem Richtungsvektor $\vec u = (2 | -1 | 2)$. Bestimme die Koordinaten eines Punkts auf $g$, der von $A$ den Abstand 9 hat. \hfill (Abitur 2019 GK)
\teil Die Gerade $g$ verläuft durch $A(0 | 4 | -1)$ mit dem Richtungsvektor $\vec u = (4 | 0 | -3)$. Bestimme die Koordinaten eines Punkts auf $g$, der von $A$ den Abstand 10 hat. \hfill (Abitur 2019 GK)
\teil Die Gerade $g$ verläuft durch $A(-2 | 1 | 3)$ mit dem Richtungsvektor $\vec u = (1 | 2 | 2)$. Bestimme die Koordinaten eines Punkts auf $g$, der von $A$ den Abstand $4{,}5$ hat. \hfill (Abitur 2019 GK)
\teil Eine Hecke verläuft geradlinig von $E(8 | 0 | 0)$ nach $F(0 | 8 | 0)$ (1 LE = 1 m). Weise nach, dass der Schattenpunkt $T(2 | 6 | 0)$ auf der Strecke $EF$ liegt. \hfill (Abitur 2018 GK)
\teil Ein Zaun verläuft geradlinig von $E(6 | -2 | 0)$ nach $F(2 | 2 | 0)$ (1 LE = 1 m). Entscheide, ob der Schattenpunkt $T(-3 | 7 | 0)$ auf dem Zaun $EF$ liegt. \hfill (Abitur 2018 GK)
\teil Ein Stahlseil ist geradlinig von $A(1 | 1 | 4)$ nach $B(7 | 4 | 1)$ gespannt (1 LE = 1 m). Weise nach, dass der Punkt $L(5 | 3 | 2)$, an dem eine Lampe hängt, auf dem Seil liegt. \hfill (Abitur 2018 GK)
\end{teile}
\end{aufgabe}

\begin{aufgabe}{Punktprobe und Punkte auf der Geraden – Prüfungsaufgaben – weiter}
\begin{teile}
\teil Weise nach, dass die Punkte $A(2 | -1 | 3)$, $B(4 | 0 | 1)$ und $C(8 | 2 | -3)$ auf einer Geraden liegen. \hfill (Abitur 2020 GK)
\teil Weise nach, dass die Punkte $A(0 | 5 | 2)$, $B(1 | 3 | 3)$ und $C(-2 | 9 | 0)$ auf einer Geraden liegen. \hfill (Abitur 2020 GK)
\teil Untersuche, ob die Punkte $A(1 | 0 | 2)$, $B(3 | 1 | 5)$ und $C(7 | 3 | 12)$ auf einer Geraden liegen. \hfill (Abitur 2020 GK)
\teil Ein Sessellift führt geradlinig von $A(1 | 20 | 2)$ nach $E(7 | 2 | 11)$; die Talstation liegt im Ursprung, die $x_1x_2$-Ebene ist waagerecht, 1 LE entspricht 5 m. Wie hoch über der Talstation ist der Lift an der Stelle, die 84 m vom Anfang $A$ entfernt ist? \hfill (Abitur 2023 GK) \\ \leerfeld[m]
\teil Eine Materialseilbahn führt geradlinig von $A(0 | 0 | 4)$ nach $E(8 | 4 | 12)$; der Ursprung liegt auf Höhe der Talstation, die $x_1x_2$-Ebene ist waagerecht, 1 LE entspricht 10 m. Wie hoch über der Talstation ist die Bahn 30 m nach dem Anfang $A$? \hfill (Abitur 2023 GK) \\ \leerfeld[m]
\teil Eine Standseilbahn fährt geradlinig von $A(3 | 1 | 0)$ nach $E(-1 | 9 | 8)$; die Talstation liegt im Ursprung, die $x_1x_2$-Ebene ist waagerecht, 1 LE entspricht 20 m. Wie hoch über der Talstation ist die Bahn 180 m nach dem Anfang $A$? \hfill (Abitur 2023 GK) \\ \leerfeld[m]
\teil Ein Prisma hat die Ecken $A(5 | 0 | 0)$ und $G(1 | 4 | 6)$. Gib eine Gleichung der Geraden durch $A$ und $G$ an und weise nach, dass der Punkt $S(3 | 2 | 5)$ nicht auf ihr liegt. \hfill (Abitur 2026 GK)
\teil Ein Quader hat die Ecken $A(0 | 6 | 1)$ und $G(4 | 0 | 4)$. Gib eine Gleichung der Geraden durch $A$ und $G$ an und weise nach, dass der Punkt $S(2 | 3 | 2)$ nicht auf ihr liegt. \hfill (Abitur 2026 GK)
\teil Ein Prisma hat die Ecken $A(2 | 0 | 0)$ und $G(-2 | 6 | 6)$. Gib eine Gleichung der Geraden durch $A$ und $G$ an und untersuche, ob der Punkt $S(0 | 3 | 3)$ auf ihr liegt. \hfill (Abitur 2026 GK)
\end{teile}
\end{aufgabe}

\begin{aufgabe}{Punkt auf einer Geraden mit vorgegebenem Abstand zu einem festen Punkt über eine quadratische Gleichung bestimmen – Prüfungsaufgaben}
\begin{teile}
\teil Von einem Mast in $F(4 | -3 | 5)$ soll ein $7{,}5$ m langes Seil zu einem Punkt der Laufschiene $g\colon \vec x = (0 | 1 | 3) + t \cdot (1 | 2 | 2)$ gespannt werden (1 LE = 1 m). Bestimme alle möglichen Befestigungspunkte auf $g$. \hfill (Abitur 2017 GK)
\end{teile}
\end{aufgabe}

% Anlauf (ohne Prüfkennung)
\begin{aufgabe}{Lage als Anhang – Anlauf}
\begin{teile}
\teil Die Richtungsvektoren der Geraden $g$ und $h$ sind Vielfache voneinander. Was ist als Nächstes zu prüfen? \kreuz{ob ein Stützpunkt von $g$ auf $h$ liegt} \\ \kreuz{ob sich die Geraden schneiden} \\ \kreuz{nichts – die Geraden sind identisch}
\teil Gegeben sind $g\colon \vec x = (2 | -1 | 3) + r \cdot (1 | 3 | 0)$ und $h\colon \vec x = (2 | -1 | 3) + t \cdot (3 | 1 | 0)$. Begründe, dass $g$ und $h$ nicht identisch sind.
\teil Ein Tunnel verläuft geradlinig durch $A(0 | 30 | 5)$ und $B(30 | 50 | 8)$, ein zweiter durch $C(10 | 0 | 2)$ und $D(70 | 40 | 8)$ (1 LE = 10 m). Prüfe, ob die beiden Tunnelachsen parallel sind.
\end{teile}
\end{aufgabe}

\begin{aufgabe}{Lage als Anhang – Prüfungsaufgaben}
\begin{teile}
\teil Die Gerade $g$ durch $P(1 | 0 | 2)$ und die Gerade $h$ durch $Q(5 | -2 | 7)$ haben beide den Richtungsvektor $(2 | -1 | 3)$. Weise nach, dass $g$ und $h$ nicht identisch sind. \hfill (Abitur 2022 GK)
\teil Die Gerade $g$ durch $P(-2 | 3 | 0)$ und die Gerade $h$ durch $Q(0 | 4 | 4)$ haben beide den Richtungsvektor $(1 | 0 | 2)$. Weise nach, dass $g$ und $h$ nicht identisch sind. \hfill (Abitur 2022 GK)
\teil Die Gerade $g$ durch $P(3 | 1 | -1)$ und die Gerade $h$ durch $Q(9 | -2 | 4)$ haben beide den Richtungsvektor $(-2 | 1 | -2)$. Weise nach, dass $g$ und $h$ nicht identisch sind. \hfill (Abitur 2022 GK)
\teil Ein Würfel hat die Ecken $A(5 | 0 | 0)$, $B(5 | 5 | 0)$, $C(0 | 5 | 0)$ und $G(0 | 5 | 5)$. Gib eine Gleichung der Geraden $h$ durch $B$ und $G$ an und begründe, dass $h$ windschief zur Geraden $AC$ ist. \hfill (Abitur 2021 GK)
\teil Eine Pyramide hat die Grundfläche $A(0 | 0 | 0)$, $B(8 | 0 | 0)$, $C(8 | 8 | 0)$, $D(0 | 8 | 0)$ und die Spitze $S(4 | 4 | 6)$. Gib eine Gleichung der Geraden $h$ durch $C$ und $S$ an und begründe, dass $h$ windschief zur Geraden $BD$ ist. \hfill (Abitur 2021 GK)
\teil Ein Quader hat die Ecken $P(3 | 0 | 0)$, $R(0 | 5 | 0)$, $T(0 | 0 | 3)$ und $U(3 | 0 | 3)$. Gib eine Gleichung der Geraden $h$ durch $R$ und $U$ an und begründe, dass $h$ windschief zur Geraden $PT$ ist. \hfill (Abitur 2021 GK)
\teil Die Punkte $A(1 | 1 | 0)$, $B(3 | 3 | 4)$, $C(0 | 0 | 6)$ und $D(2 | 2 | 4)$ liegen in der Ebene $E\colon x_1 - x_2 = 0$. Die Gerade $g$ verläuft durch $A$ und $B$, die Gerade $h$ durch $C$ und $D$. Begründe, ohne den Schnittpunkt zu berechnen, dass $g$ und $h$ sich schneiden müssen. \hfill (Abitur 2018 GK)
\teil Zwei Balken liegen auf einer waagerechten Decke: der erste von $A(0 | 1 | 2)$ nach $B(4 | 3 | 2)$, der zweite von $C(1 | 5 | 2)$ nach $D(3 | -1 | 2)$ (1 LE = 1 m). Begründe, ohne den Schnittpunkt zu berechnen, dass sich die Geraden der beiden Balken schneiden. \hfill (Abitur 2018 GK)
\teil Die Punkte $A(0 | 0 | 6)$, $B(3 | 4 | 0)$, $C(1 | 5 | 4)$ und $D(2 | -1 | 2)$ liegen in der Ebene $E\colon 2x_1 + x_3 = 6$. Begründe, ohne den Schnittpunkt zu berechnen, dass sich die Gerade durch $A$ und $B$ und die Gerade durch $C$ und $D$ schneiden. \hfill (Abitur 2018 GK)
\teil Gegeben sind $g\colon \vec x = (1 | 3 | 1) + r \cdot (2 | 0 | 1)$ und die dazu echt parallele Gerade $p\colon \vec x = (1 | 5 | 1) + r \cdot (2 | 0 | 1)$. Entscheide, ob jede Gerade, die $g$ senkrecht schneidet, auch $p$ schneidet, und begründe. \hfill (Abitur 2019 GK)
\teil Gegeben sind $g\colon \vec x = (0 | 2 | 3) + r \cdot (1 | 1 | 0)$ und die dazu echt parallele Gerade $p\colon \vec x = (0 | 2 | 5) + r \cdot (1 | 1 | 0)$. Entscheide, ob jede Gerade, die $g$ senkrecht schneidet, auch $p$ schneidet, und begründe. \hfill (Abitur 2019 GK)
\end{teile}
\end{aufgabe}

% Anlauf (ohne Prüfkennung)
\begin{aufgabe}{Sachgeraden – Anlauf}
\begin{teile}
\teil Ein Seil ist zwischen zwei Masten gespannt. Wird es im Modell durch eine Gerade oder durch eine Strecke beschrieben? \kreuz{Gerade} \\ \kreuz{Strecke}
\teil Ein Seil ist an zwei Masten in $P(0 | 0 | 8)$ und $Q(6 | 8 | 5)$ befestigt (1 LE = 1 m). Gib eine Gleichung der Seilstrecke an.
\end{teile}
\end{aufgabe}

\begin{aufgabe}{Sachgeraden – Prüfungsaufgaben}
\begin{teile}
\teil Ein unteres Seil ist geradlinig von $A(0 | 0 | 6)$ nach $B(4 | 8 | 4)$ gespannt, ein oberes verläuft entlang $h\colon \vec x = (-2 | 2 | 10) + t \cdot (3 | 0 | 0)$ (1 LE = 1 m). Ein Punkt des oberen Seils liegt senkrecht über einem Punkt des unteren. Berechne den Abstand dieser beiden Punkte. \hfill (Abitur 2022 GK) \\ \leerfeld[m]
\teil Eine Stromleitung verläuft entlang $g\colon \vec x = (2 | 0 | 3) + s \cdot (2 | 6 | 1)$, ein Tragseil darüber entlang $h\colon \vec x = (0 | 3 | 8) + t \cdot (2 | 0 | -1)$ (1 LE = 1 m). Ein Punkt von $h$ liegt senkrecht über einem Punkt von $g$. Berechne den Abstand dieser beiden Punkte. \hfill (Abitur 2022 GK) \\ \leerfeld[m]
\teil Eine Wasserleitung verläuft waagerecht entlang $g\colon \vec x = (1 | 1 | 2) + s \cdot (2 | 2 | 0)$, ein Stahlseil entlang $h\colon \vec x = (0 | 6 | 9) + t \cdot (1 | -1 | -1)$ (1 LE = 1 m). Ein Punkt des Seils liegt senkrecht über einem Punkt der Leitung. Berechne den Abstand dieser beiden Punkte. \hfill (Abitur 2022 GK) \\ \leerfeld[m]
\teil Ein Seil ist geradlinig zwischen zwei Masten gespannt; die Fußpunkte sind $F_1(0 | 0 | 0)$ und $F_2(6 | 8 | 0)$, befestigt ist es in 9 m und in $7{,}5$ m Höhe (1 LE = 1 m). Berechne die Neigung des Seils in Prozent. \hfill (Abitur 2022 GK) \\ \leerfeld \%
\teil Ein Seil ist geradlinig zwischen zwei Masten gespannt; die Fußpunkte sind $F_1(0 | 0 | 0)$ und $F_2(12 | 5 | 0)$, befestigt ist es in 10 m und in $7{,}4$ m Höhe (1 LE = 1 m). Berechne die Neigung des Seils in Prozent. \hfill (Abitur 2022 GK) \\ \leerfeld \%
\teil Ein Seil ist geradlinig zwischen zwei Masten gespannt; die Fußpunkte sind $F_1(3 | 1 | 0)$ und $F_2(11 | 16 | 0)$, befestigt ist es in 8 m und in $6{,}3$ m Höhe (1 LE = 1 m). Berechne die Neigung des Seils in Prozent. \hfill (Abitur 2022 GK) \\ \leerfeld \%
\teil Ein Sonnensegel wird mit einem geradlinig gespannten Seil befestigt. Das Seil beginnt am ersten Mast im Punkt $(0 | 0 | 3)$ und endet am zweiten Mast, der senkrecht durch $M(6 | 8 | 0)$ steht. Unterwegs liegt es auf der Kante $UV$ eines waagerechten Dachs in 4 m Höhe auf, mit $U(4 | 3 | 4)$ und $V(1 | 6 | 4)$; das Verhältnis, in dem der Auflagepunkt die Kante teilt, ist bekannt (1 LE = 1 m). Beschreibe ein Verfahren, mit dem sich der Abstand des Befestigungspunkts am zweiten Mast vom Dach berechnen lässt. \hfill (Abitur 2018 GK)
\teil Ein Seil beginnt an einer Plattform in $S(1 | 0 | 2)$, liegt auf der Oberkante $KL$ einer Mauer in 3 m Höhe auf, mit $K(8 | 2 | 3)$ und $L(4 | 6 | 3)$, und ist an einem senkrechten Pfosten durch $P(15 | 10 | 0)$ befestigt. Das Verhältnis, in dem der Auflagepunkt die Kante $KL$ teilt, ist bekannt (1 LE = 1 m). Beschreibe ein Verfahren, mit dem sich berechnen lässt, wie hoch über der Mauerkrone das Seil am Pfosten befestigt ist. \hfill (Abitur 2018 GK)
\teil Eine Kamera gleitet auf einer Schiene von $A(1 | 0 | 6)$ nach $B(5 | 2 | 4)$ und filmt stets in Richtung $(2 | 1 | -3)$ (1 LE = 1 m). Um zu prüfen, ob sie den Punkt $F(8 | 3 | 0{,}5)$ filmen kann, wird der Ansatz $\vec{OK_t} + r \cdot (2 | 1 | -3) = (8 | 3 | 0{,}5)$ mit $\vec{OK_t} = (1 | 0 | 6) + t \cdot (4 | 2 | -2)$, $0 \le t \le 1$, aufgestellt. Erläutere die geometrischen Sachverhalte des Ansatzes und deute ihn im Sachzusammenhang. \hfill (Abitur 2025 GK)
\teil Ein Scheinwerfer fährt auf einer Schiene von $P(0 | 2 | 7)$ nach $Q(6 | 2 | 5)$ und strahlt stets in Richtung $(1 | 1 | -2)$ (1 LE = 1 m). Um zu prüfen, ob der Punkt $Z(7 | 6 | 0)$ auf der Bühne beleuchtet werden kann, wird der Ansatz $(0 | 2 | 7) + t \cdot (6 | 0 | -2) + r \cdot (1 | 1 | -2) = (7 | 6 | 0)$ mit $t \in [0; 1]$ aufgestellt. Erläutere die Bestandteile des Ansatzes und deute ihn im Sachzusammenhang. \hfill (Abitur 2025 GK)
\end{teile}
\end{aufgabe}

\begin{aufgabe}{Horizontalabstand aus vorgegebener Neigung in Prozent und Höhendifferenz berechnen – Prüfungsaufgaben}
\begin{teile}
\teil Ein Steg beginnt am rechten Ufer in $0{,}9$ m Höhe über dem Wasserspiegel und endet links auf der Uferzone in $0{,}3$ m Höhe; er soll eine Steigung von 4\,\% haben. Im Querschnitt liegt der Anfang bei $x = 0$, die linke Uferzone beginnt bei $x = -10$ (Angaben in m). Wie weit liegt das linke Ende des Stegs vom Rand der Uferzone entfernt? \hfill (Abitur 2017 GK) \\ \leerfeld[m]
\teil Eine Rampe mit 6\,\% Steigung soll vor einer Haustür 45 cm Höhenunterschied überwinden. Zwischen Tür und Gehweg liegen 4 m Vorgarten. Ragt die Rampe in den Gehweg hinein – und wenn ja, wie weit? \hfill (Abitur 2017 GK) \\ \leerfeld[m]
\teil Eine Rampe beginnt in $B(0 | 0 | 12)$, führt in Richtung der $x_1$-Achse mit 5\,\% Gefälle geradlinig abwärts und endet auf einer Straße in 3 m Höhe (1 LE = 1 m). Gib den Endpunkt der Rampe an. \hfill (Abitur 2017 GK)
\end{teile}
\end{aufgabe}

\clearpage
% Ebenen – bank/ebenen/ (zusammenbau v0.4, Heft)
\einheitenkopf[t4]{Ebenen}
\setcounter{aufgabe}{36}

% Anlauf (ohne Prüfkennung)
\begin{aufgabe}{Parameterform – Anlauf}
\begin{teile}
\teil Eine Ebene soll durch $A(1 | 3 | 0)$, $B(2 | 0 | 4)$ und $C(5 | 1 | 1)$ gehen. Was ist gegeben? \\ \kreuz{drei Punkte} \\ \kreuz{ein Punkt und zwei Richtungen} \\ \kreuz{eine Gerade und ein Punkt} \\ \kreuz{zwei Geraden}
\teil Punkt $(2 | 1 | 3)$, Spannvektoren $(1 | 0 | 4)$ und $(0 | 3 | 1)$: Parametergleichung der Ebene $E$?
\teil $A(2 | 1 | -1)$, $B(4 | 2 | 0)$, $C(0 | 3 | 1)$: Parametergleichung der Ebene durch $A$, $B$, $C$? Prüfe, dass die Spannvektoren nicht kollinear sind.
\end{teile}
\end{aufgabe}

\begin{aufgabe}{Parameterform – Prüfungsaufgaben}
\begin{teile}
\teil Ein Zeltdach hat vier rautenförmige Dachflächen; eine davon hat die Ecken $C(6 | 6 | 0)$, $F(6 | 3 | 3)$, $G(3 | 6 | 3)$ und die Spitze $S$ (1 LE = 1 m). Gib eine Parametergleichung der Ebene $L$ durch $C$, $F$ und $G$ an und weise nach, dass $S(3 | 3 | 6)$ in $L$ liegt. \hfill (Abitur 2022 GK)
\teil Gegeben sind $A(2 | -1 | 1)$, $B(1 | 3 | 4)$ und der Punkt $C$ mit $\overrightarrow{OC} = 3 \cdot \overrightarrow{OA}$; $O$ ist der Koordinatenursprung. Begründe, dass es genau eine Ebene gibt, die $O$, $A$, $B$ und $C$ enthält. \hfill (Abitur 2018 GK)
\teil Gegeben sind $A(-1 | 2 | 2)$, $B(4 | 0 | 1)$ und der Punkt $C$ mit $\overrightarrow{OC} = -2 \cdot \overrightarrow{OA}$; $O$ ist der Koordinatenursprung. Begründe, dass es genau eine Ebene gibt, die $O$, $A$, $B$ und $C$ enthält. \hfill (Abitur 2018 GK)
\teil Das Pultdach eines Carports ist das Rechteck mit den Ecken $E(1 | 0 | 3)$, $F(7 | 0 | 3)$, $G(7 | 5 | 6)$ und $H(1 | 5 | 6)$ (1 LE = 1 m). Gib eine Parametergleichung der Dachebene $L$ an und prüfe, ob der Befestigungspunkt $P(4 | 2{,}5 | 4{,}5)$ einer Lampe in $L$ liegt. \hfill (Abitur 2022 GK)
\teil Die Dachfläche eines Gartenhauses ist das Rechteck mit den Ecken $A(0 | 0 | 3)$, $B(6 | 0 | 3)$, $C(6 | 4 | 5)$ und $D(0 | 4 | 5)$ (1 LE = 1 m). Gib eine Parametergleichung der Dachebene $L$ an und prüfe, ob der Punkt $P(3 | 3 | 4)$ in $L$ liegt. \hfill (Abitur 2022 GK)
\end{teile}
\end{aufgabe}

% Anlauf (ohne Prüfkennung)
\begin{aufgabe}{Koordinatengleichung bestimmen – Anlauf}
\begin{teile}
\teil Gesucht ist eine Gleichung der Form $ax + by + cz = d$ für die Ebene durch drei gegebene Punkte. Welche Form ist verlangt? \\ \kreuz{Parameterform} \\ \kreuz{Koordinatenform} \\ \kreuz{Normalenform}
\teil Spannvektoren $(1 | 0 | 3)$ und $(0 | 1 | 2)$: ein Normalenvektor $\vec n$ der Ebene? $\vec n = ($ \leerfeld | \leerfeld | \leerfeld $)$
\teil Eine Ebene schneidet die Achsen in $(6 | 0 | 0)$, $(0 | 3 | 0)$ und $(0 | 0 | 2)$. Koordinatengleichung?
\end{teile}
\end{aufgabe}

\begin{aufgabe}{Koordinatengleichung bestimmen – Prüfungsaufgaben}
\begin{teile}
\teil Ein dreieckiges Sonnensegel ist an drei Masten in $A(3 | 2 | 2)$, $B(3 | 5 | 2)$ und $C(0 | 3 | 4)$ befestigt (1 LE = 1 m). Die Ebene des Segels hat eine Gleichung der Form $2x + 3z = d$. Bestimme $d$. \hfill (Abitur 2020 GK) \\ $d = $ \leerfeld
\teil Ein Dachaufbau hat die Form einer Pyramide mit der Grundfläche $A(1 | 1 | 0)$, $B(4 | 1 | 0)$, $C(4 | 4 | 0)$, $D(1 | 7 | 0)$ und der Spitze $S(1 | 1 | 5)$ (1 LE = 1 m). Die Seitenfläche $CDS$ liegt in der Ebene $E$. Bestimme eine Gleichung von $E$ in Koordinatenform (zur Kontrolle: $5x + 5y + 6z = 40$). \hfill (Abitur 2024 GK)
\teil Die Seitenfläche $BCS$ eines pyramidenförmigen Kristalls hat die Ecken $B(3 | 1 | 3)$, $C(2 | 4 | 3)$ und $S(1 | 1 | 6)$ (1 LE = 1 mm). Bestimme eine Koordinatengleichung der Ebene $E$, in der $BCS$ liegt (zur Kontrolle: $3x + y + 2z = 16$). \hfill (Abitur 2025 GK)
\teil Ein Holzkörper ist eine Pyramide über dem Rechteck $ABCD$ mit $B(6 | 0 | 0)$ und $C(6 | 8 | 0)$; seine Spitze $E(0 | 8 | 4)$ liegt senkrecht über $D$ (1 LE = 1 cm). Bestimme eine Koordinatengleichung der Ebene $L$ durch $B$, $C$ und $E$. \hfill (Abitur 2021 GK)
\teil Bestimme eine Koordinatengleichung von $E\colon \vec x = (4 | 1 | 0) + r \cdot (-2 | 4 | 0) + s \cdot (-2 | 0 | 3)$, $r, s \in \mathbb{R}$ (zur Kontrolle: $6x + 3y + 4z = 27$). \hfill (Abitur 2020 GK)
\teil Ein Sonnensegel $W$ ist in $M(3 | 3 | 2)$, $F(0 | 0 | 3)$ und $S(6 | 0 | 0)$ befestigt (1 LE = 1 m). Bestimme eine Koordinatengleichung von $W$ mit dem Ansatz $ax + by + cz = 18$ (zur Kontrolle: $3x - y + 6z = 18$). \hfill (Abitur 2024 GK)
\teil Eine dreieckige Plane über einer Lagerhalle hat die Ecken $U(1 | 1 | 5)$, $V(3 | 9 | 5)$ und $W(9 | 1 | 3)$ (1 LE = 1 m). Gib die Ebene $E$ der Plane in Parameter- und in Koordinatenform an (zur Kontrolle: $4x - y + 16z = 83$). \hfill (Abitur 2020 GK)
\teil Die seitlichen Streben einer Brücke liegen in der Ebene $E$ durch $A(1 | 0 | 6)$, $B(3 | 40 | 6)$ und $C(1{,}2 | 4 | 8)$; die Fahrbahn liegt in der $xy$-Ebene (1 LE = 1 m). Bestimme einen Normalenvektor mit dem Vektorprodukt und gib $E$ in Parameter- und in Koordinatenform an. \hfill (Abitur 2018 GK)
\end{teile}
\end{aufgabe}

% Anlauf (ohne Prüfkennung)
\begin{aufgabe}{Normalenvektor und Nachweise – Anlauf}
\begin{teile}
\teil Gegeben ist $E\colon \vec x = \vec p + r \cdot \vec u + s \cdot \vec v$ mit Zahlen für die Vektoren. In welcher Form ist $E$ gegeben? \\ \kreuz{Parameterform} \\ \kreuz{Koordinatenform} \\ \kreuz{Normalenform}
\teil Gib zwei Normalenvektoren von $E\colon 4x - 2y + z = 7$ an.
\teil Gib $E\colon \left(\vec x - (1 | -2 | 0)\right) \cdot (2 | 1 | 5) = 0$ in Koordinatenform an.
\end{teile}
\end{aufgabe}

\begin{aufgabe}{Normalenvektor und Nachweise – Prüfungsaufgaben}
\begin{teile}
\teil Gegeben ist $F\colon \left(\vec x - (2 | 5 | 1)\right) \cdot (-3 | 1 | 4) = 0$. Gib $F$ in Koordinatenform an. \hfill (Abitur 2021 GK)
\teil Gegeben sind $E(5 | 2 | 3)$, $F(1 | 4 | 3)$ und $S(1 | 1 | 5)$. Bestimme rechnerisch einen Normalenvektor der Ebene $EFS$ und prüfe ihn. \hfill (Abitur 2025 GK)
\end{teile}
\end{aufgabe}

% Anlauf (ohne Prüfkennung)
\begin{aufgabe}{Besondere Lagen – Anlauf}
\begin{teile}
\teil $E\colon 4x + 5z = 30$. Welche Variable fehlt? \\ \kreuz{$x$} \\ \kreuz{$y$} \\ \kreuz{$z$} \\ \kreuz{keine}
\teil Gleichung der $xz$-Ebene?
\teil Zu welcher Achse ist $E\colon 5x + z = 5$ parallel? Begründe.
\end{teile}
\end{aufgabe}

\begin{aufgabe}{Besondere Lagen – Prüfungsaufgaben}
\begin{teile}
\teil Welche besondere Lage hat $E\colon 4x - 5y = 0$? Beschreibe sie. \hfill (Abitur 2022 GK)
\teil Eine Kletterwand ist ein ebenes Viereck mit $A(2 | 6 | 1)$, $B(6 | 4 | 1)$, $C(5 | 4{,}5 | 5)$ und $D(1 | 6{,}5 | 5)$; die $xy$-Ebene ist der Boden (1 LE = 1 m). Weise nach, dass die Wand vertikal steht. \hfill (Abitur 2022 GK)
\end{teile}
\end{aufgabe}

% Anlauf (ohne Prüfkennung)
\begin{aufgabe}{Parallele Ebenen – Anlauf}
\begin{teile}
\teil $E\colon x + 2y - z = 3$ und $F\colon 2x + 4y - 2z = 6$. Wie liegen sie zueinander? \\ \kreuz{parallel} \\ \kreuz{identisch} \\ \kreuz{schneidend}
\teil Ebene $F$ parallel zu $E\colon 3x + y - 2z = 4$ durch $P(2 | 1 | 1)$?
\teil $F$ ist parallel zu $E\colon 2x - y + 2z = 4$ und geht durch $P(1 | 1 | 3)$. Liegt $Q(3 | 1 | 1)$ in $F$?
\end{teile}
\end{aufgabe}

\begin{aufgabe}{Parallele Ebenen – Prüfungsaufgaben}
\begin{teile}
\teil Ein Würfel mit der Kantenlänge 6 hat eine Ecke im Ursprung und liegt im ersten Oktanten. Die Punkte $K(0 | 6 | 3)$ und $L(1 | 0 | 6)$ liegen auf Kanten. $K$ liegt in einer Ebene $T$, die parallel zu $S\colon 3x + 2y + 3z = 0$ ist. Liegt auch $L$ in $T$? \hfill (Abitur 2019 GK)
\teil Bei einem Körper liegen die Ecken $A$, $B$, $C$ in der Ebene $L_1\colon 2x + y + 2z = 8$, die Ecken $D$, $E$, $F$ in $L_2\colon 4x + 2y + 4z = 14$. Begründe, dass $L_1$ und $L_2$ keine gemeinsamen Punkte haben. \hfill (Abitur 2023 GK)
\teil Eine Halle hat die quadratische Grundfläche $ABCD$ mit 6 m Seitenlänge in der $xy$-Ebene, $A$ im Ursprung, $B$ auf der $x$-Achse, $D$ auf der $y$-Achse (1 LE = 1 m). Ein Zuschauerboden liegt in der Ebene $E_B$, die parallel zu $P\colon x - y + 6z = 30$ durch $B$ verläuft. Gib $E_B$ an und bestimme die größte Höhe des Bodens über der Grundfläche. \hfill (Abitur 2020 GK)
\teil Ein dreiseitiges Prisma hat die Grundfläche $ABC$ in der Ebene $L\colon x + 2y + 2z = 3$ und die Deckfläche $DEF$; die Seitenkante $BE$ verbindet $B(1 | 1 | 0)$ mit $E(4 | 4 | 6)$. Eine zu $L$ parallele Ebene $M$ teilt das Prisma so, dass der Teilkörper, der $E$ enthält, doppelt so groß ist wie der andere. Bestimme eine Gleichung von $M$. \hfill (Abitur 2020 GK)
\teil Ein Prisma hat die Grundfläche in der Ebene $L\colon 2x - y + 2z = 1$; eine Seitenkante verbindet $A(1 | 1 | 0)$ in der Grundfläche mit $D(5 | 5 | 4)$ in der Deckfläche. Eine zu $L$ parallele Ebene $M$ teilt das Prisma so, dass der Teilkörper, der $D$ enthält, dreimal so groß ist wie der andere. Bestimme eine Gleichung von $M$. \hfill (Abitur 2020 GK)
\end{teile}
\end{aufgabe}

% Anlauf (ohne Prüfkennung)
\begin{aufgabe}{Parallele Ebenen über die Normalenvektoren auswählen und über die Konstante der Koordinatengleichung der Höhe nach zuordnen – Anlauf}
\begin{teile}
\teil Schräge Regalböden liegen in Ebenen ($z$ zeigt nach oben): $E_1\colon x + 3z = 6$, $E_2\colon x + 3z = 9$, $E_3\colon 2x + 6z = 30$, $E_4\colon x + 2z = 9$. Welche sind parallel zu $E_1$, und in welcher Reihenfolge liegen sie von unten nach oben?
\end{teile}
\end{aufgabe}

\begin{aufgabe}{Parallele Ebenen über die Normalenvektoren auswählen und über die Konstante der Koordinatengleichung der Höhe nach zuordnen – Prüfungsaufgaben}
\begin{teile}
\teil Ein Turm hat drei Dachetagen; die Dachflächen der Südseite sind parallel ($z$ zeigt nach oben, 1 LE = 1 m). Die untere liegt in $E\colon 2x + 5z = 40$, die mittlere und die obere in zwei der Ebenen I: $2x + 4z = 40$, II: $2x + 5z = 30$, III: $2x + 5z = 50$, IV: $4x + 5z = 40$, V: $2x + 5z = 60$. Ordne zu und begründe. \hfill (Abitur 2017 GK)
\end{teile}
\end{aufgabe}

\clearpage
% Lagebeziehungen – bank/lagebeziehungen/ (zusammenbau v0.4, Heft)
\einheitenkopf[t5]{Lagebeziehungen}
\setcounter{aufgabe}{48}

% Anlauf (ohne Prüfkennung)
\begin{aufgabe}{Punktprobe und Seitenlage – Anlauf}
\begin{teile}
\teil $E: 2x - y + z = 3$. Die Koordinaten von $P$ eingesetzt ergeben links $6$. Wo liegt $P$? \\ \kreuz{in $E$} \\ \kreuz{auf der Seite mit größerem Wert} \\ \kreuz{auf der Seite mit kleinerem Wert}
\teil Liegt $P(3 | 1 | 4)$ in $E: 2x - y + z = 7$?
\teil $E_1: 2x - y + 2z = 1$ und $E_2: 2x - y + 2z = 7$ sind parallel. Liegt $P(1 | 1 | 2)$ zwischen den beiden Ebenen? Begründe.
\end{teile}
\end{aufgabe}

\begin{aufgabe}{Punktprobe und Seitenlage – Prüfungsaufgaben}
\begin{teile}
\teil $E: \vec{x} = (2 | 0 | 1) + t \cdot (1 | 1 | 0) + s \cdot (0 | 1 | 2)$. Weise nach, dass $P(4 | 5 | 7)$ in $E$ liegt. \hfill (Abitur 2021 GK)
\teil $E: \vec{x} = (1 | 1 | 0) + t \cdot (2 | 0 | 1) + s \cdot (0 | 1 | 2)$. Prüfe, ob $P(5 | 3 | 5)$ in $E$ liegt. \hfill (Abitur 2021 GK)
\teil Eine Lärmschutzwand steht auf der Fahrbahn (xy-Ebene) und liegt in der Ebene $W: 5x - 2y = 0$ (1 LE = 1 m). Weise nach, dass der Punkt $K(2 | 5 | 3)$ in $W$ liegt und dass $W$ senkrecht auf der Fahrbahn steht. \hfill (Abitur 2018 GK)
\teil $P(1 | 0 | -2)$, $Q(3 | 4 | 0)$ und $E: x + 2y + z = 11$. Weise nach, dass $Q$ in $E$ liegt und dass $\vec{PQ}$ ein Normalenvektor von $E$ ist. \hfill (Abitur 2025 GK)
\teil $E: 2x + y - 2z = 3$; die Gerade $g$ geht durch $P(1 | 3 | 1)$ und $Q(5 | 5 | -3)$. Weise nach, dass $P$ in $E$ liegt, und begründe, dass $g$ senkrecht auf $E$ steht. \hfill (Abitur 2026 GK)
\teil Ein Quader hat die Kante $AE$ mit $A(5 | 0 | 0)$ und $E(5 | 0 | 3)$; die Ebene $W: x + 2y + 2z = 9$ ist gegeben. Ein Lösungsweg lautet: (1) $P(5 | 0 | r)$ mit $0 \le r \le 3$; (2) $5 + 2 \cdot 0 + 2r = 9$. Formuliere eine passende Aufgabenstellung und führe sie zu Ende. \hfill (Abitur 2024 GK)
\teil Ein Körper hat die Kante $BC$ mit $B(2 | 6 | 0)$ und $C(2 | 6 | 4)$; die Ebene $W: 3x - y + z = 1$ ist gegeben. Ein Lösungsweg lautet: (1) $P(2 | 6 | r)$ mit $0 \le r \le 4$; (2) $3 \cdot 2 - 6 + r = 1$. Formuliere eine passende Aufgabenstellung und führe sie zu Ende. \hfill (Abitur 2024 GK)
\teil Ein Quader hat die Kante $FG$ mit $F(0 | 3 | 2)$ und $G(6 | 3 | 2)$; die Ebene $W: x - y + 2z = 5$ ist gegeben. Ein Lösungsweg lautet: (1) $P(r | 3 | 2)$ mit $0 \le r \le 6$; (2) $r - 3 + 2 \cdot 2 = 5$. Formuliere eine passende Aufgabenstellung und führe sie zu Ende. \hfill (Abitur 2024 GK)
\teil Ein Körper ist der Teil des ersten Oktanten zwischen den parallelen Ebenen $L_1: x + 2y + z = 2$ und $L_2: x + 2y + z = 6$. Begründe, dass $T(1 | 0{,}5 | 1)$ im Inneren des Körpers liegt. \hfill (Abitur 2023 GK)
\teil Ein Haus ist vereinfacht der Quader mit $0 \le x \le 8$, $0 \le y \le 4$ und $0 \le z \le 5$ (1 LE = 1 m). Untersuche für die Ebenen $S: x - 2y = 0$ und $T: x + 2y = 0$, ob sie durch das Innere des Hauses verlaufen. \hfill (Abitur 2019 GK)
\end{teile}
\end{aufgabe}

% Anlauf (ohne Prüfkennung)
\begin{aufgabe}{Parameter aus der Lagebedingung – Anlauf}
\begin{teile}
\teil Für welche Zahl $k$ liegt $P(1 | 2 | 0)$ in der Ebene $E_k: kx + 3y + z = 8$? Wer gegen wen? \\ \kreuz{Punkt gegen Ebene} \\ \kreuz{Gerade gegen Ebene} \\ \kreuz{Punkt gegen Gerade}
\teil $E_k: kx + 2y - z = 5$ enthält $P(2 | 1 | 3)$. Bestimme $k$. $k = $ \leerfeld
\teil Eine Ebene der Form $r \cdot x + s \cdot z = 0$ enthält den Punkt $P(4 | 1 | -2)$. Gib passende Werte für $r$ und $s$ an. $r = $ \leerfeld , $s = $ \leerfeld
\end{teile}
\end{aufgabe}

\begin{aufgabe}{Parameter aus der Lagebedingung – Prüfungsaufgaben}
\begin{teile}
\teil Das Dreieck mit $A(3 | 0 | 0)$, $B(0 | 6 | 0)$ und $C(0 | 0 | 2)$ liegt in einer Ebene der Form $2x + y + kz = 6$. Bestimme $k$. \hfill (Abitur 2023 GK) \\ $k = $ \leerfeld
\teil Die Ebene $L$ enthält die x-Achse und den Punkt $G(3 | 2 | 5)$; ihre Gleichung hat die Form $r \cdot y + s \cdot z = 0$. Gib passende Werte für $r$ und $s$ an. \hfill (Abitur 2023 GK) \\ $r = $ \leerfeld , $s = $ \leerfeld
\end{teile}
\end{aufgabe}

% Anlauf (ohne Prüfkennung)
\begin{aufgabe}{Gerade und Ebene – Anlauf}
\begin{teile}
\teil Der allgemeine Punkt von $g$ in die Gleichung von $E$ eingesetzt ergibt $4 + 3t - 3t = 4$. Wie liegt $g$ zu $E$? \\ \kreuz{liegt in der Ebene} \\ \kreuz{echt parallel} \\ \kreuz{genau ein Schnittpunkt}
\teil Weise nach, dass $g: \vec{x} = (1 | 2 | 0) + t \cdot (1 | -1 | 1)$ in $E: x + 2y + z = 5$ liegt.
\teil $g: \vec{x} = (1 | 0 | 2) + t \cdot (2 | 1 | 1)$ und $E: x - 2y + z = 1$. Entscheide mit der Merkregel, wie $g$ zu $E$ liegt.
\end{teile}
\end{aufgabe}

\begin{aufgabe}{Gerade und Ebene – Prüfungsaufgaben}
\begin{teile}
\teil Weise nach, dass die Gerade $s: \vec{x} = (2 | -1 | 3) + r \cdot (0 | 6 | -3)$ in $E: x + y + 2z = 7$ liegt. \hfill (Abitur 2021 GK)
\teil Eine Drohne fliegt auf $h: \vec{x} = (100 | 50 | 40) + s \cdot (30 | 10 | 1)$. Die Unterseite einer Wolkendecke liegt in $E: x - 30z = -800$ (1 LE = 1 m). Weise nach, dass die Flugbahn parallel zu $E$ verläuft, und prüfe, ob sie in $E$ liegt. \hfill (Abitur 2017 GK)
\teil $P(1 | 3 | 8)$ und $Q_t(5 | t | 6)$ mit reellem $t$. Gibt es ein $t$, für das die Gerade $PQ_t$ parallel zur xy-Ebene verläuft? Begründe. \hfill (Abitur 2024 GK)
\teil $A(3 | 1 | 2)$ und $B(5 | -4 | 2)$. Begründe, dass die Gerade $AB$ parallel zur xy-Ebene verläuft. \hfill (Abitur 2017 GK)
\end{teile}
\end{aufgabe}

% Anlauf (ohne Prüfkennung)
\begin{aufgabe}{Lagebefunde im Sachzusammenhang – Anlauf}
\begin{teile}
\teil Frage: Liegt der Schatten der Lampe auf der Wand? Rechnung: Durchstoßpunkt der Lichtgeraden mit der Wandebene $(2 | 0 | 1{,}8)$. Ist die Frage damit beantwortet? \\ \kreuz{ja} \\ \kreuz{nein – ein Bereich ist noch zu prüfen}
\teil Eine Wand liegt in der Ebene $y = 6$ mit $0 \le x \le 6$ und $0 \le z \le 3$. Eine Lampe in $L(2 | 0 | 4)$ beleuchtet die Spitze $P(3 | 2 | 3)$ eines Pfostens. Liegt der Schatten von $P$ auf der Wand?
\teil Eine Terrasse ist das Rechteck mit $0 \le x \le 4$ und $0 \le y \le 6$ in der xy-Ebene. Ein Lösungsschritt lautet: $(x | y | 0) = (2 | 1 | 3) + \lambda \cdot (1 | 1 | -1)$ liefert $\lambda = 3$ und $(5 | 4 | 0)$. Erläutere den Schritt und ziehe eine Folgerung für den Schatten der Ecke $(2 | 1 | 3)$ eines Sonnensegels.
\end{teile}
\end{aufgabe}

\begin{aufgabe}{Lagebefunde im Sachzusammenhang – Prüfungsaufgaben}
\begin{teile}
\teil Sonnenlicht fällt in Richtung $(-1 | -2 | -1)$; eine Hauswand liegt in der xz-Ebene mit $0 \le x \le 6$ und $0 \le z \le 3$. Schritt I: $(x | 0 | z) = (4 | 2 | 3) + \lambda \cdot (-1 | -2 | -1)$ liefert $\lambda = 1$ und $(3 | 0 | 2)$. Schritt II: $0 < 3 < 6$ und $0 < 2 < 3$. Erläutere beide Schritte im Sachzusammenhang. \hfill (Abitur 2024 GK)
\teil Eine Wand liegt in der xz-Ebene mit $0 \le x \le 5$ und $0 \le z \le 4$; $(1 | 2 | 4)$ ist eine Ecke der Unterkante eines Rollos. Ein Lösungsschritt lautet: $(x | 0 | z) = (1 | 2 | 4) + \mu \cdot (1 | -2 | -1)$ liefert $\mu = 1$ und $(2 | 0 | 3)$. Erläutere die Bedeutung dieses Schritts und ziehe eine Folgerung. \hfill (Abitur 2023 GK)
\end{teile}
\end{aufgabe}

\clearpage
% Schnittmengen – bank/schnittmengen/ (zusammenbau v0.4, Heft)
\einheitenkopf[t6]{Schnittmengen}
\setcounter{aufgabe}{56}

% Anlauf (ohne Prüfkennung)
\begin{aufgabe}{Schnittpunkt von Gerade und Ebene – Anlauf}
\begin{teile}
\teil Gegeben sind eine Gerade in Parameterform und eine Ebene in Koordinatenform. Welcher Weg führt zum Schnittpunkt? \\ \kreuz{einsetzen} \\ \kreuz{gleichsetzen}
\teil Bestimme den Schnittpunkt der Geraden $\vec{x} = (-1|2|1) + t \cdot (1|0|1)$ mit der Ebene $x + y + z = 6$. S( \leerfeld | \leerfeld | \leerfeld )
\end{teile}
\end{aufgabe}

\begin{aufgabe}{Schnittpunkt von Gerade und Ebene – Prüfungsaufgaben}
\begin{teile}
\teil Ein Körper hat die Spitze $P(0|0|3)$; $M(3|1|0)$ ist der Mittelpunkt einer Grundkante. Die Gerade durch $P$ und $M$ schneidet die Ebene $L$: $x + 2y + z = 4$ im Punkt $Q$. Bestimme die Koordinaten von $Q$. \hfill (Abitur 2023 GK) \\ Q( \leerfeld | \leerfeld | \leerfeld )
\teil Aus einem Versorgungsstollen führt von $L(4|6|2)$ ein Rohr senkrecht nach oben durch den Berg und ragt 1,5 m aus ihm heraus. Die Bergoberfläche liegt in der Ebene $F$: $2x + 3y + 5z = 51$; die $xy$-Ebene liegt auf Meereshöhe, eine Längeneinheit entspricht 50 m. Berechne die Länge des Rohrs. \hfill (Abitur 2019 GK) \\ \leerfeld[m]
\teil Ein Sonnensegel hat die Ecken $A(1|2|2)$, $B(1|6|2)$ und $C(-2|4|4)$; der Boden ist die $xy$-Ebene, eine Längeneinheit entspricht 1 m. Bei parallelem Sonnenlicht liegen die Schatten von $A$ und $B$ in $A'(-1|3|0)$ und $B'(-1|7|0)$. Bestimme den Schatten $C'$ von $C$. \hfill (Abitur 2020 GK) \\ C'( \leerfeld | \leerfeld | \leerfeld )
\teil Durch die Glasfassade eines Hauses fällt Sonnenlicht in Richtung $(2|-1|-2)$; eine Längeneinheit entspricht 1 m. Der Boden des Dachgeschosses liegt in der Ebene $z = 6$, der obere Eckpunkt der Fassade ist $L(0|5|9)$. Bestimme den Punkt, in dem das Licht durch $L$ auf den Boden des Dachgeschosses trifft. \hfill (Abitur 2019 GK) \\ ( \leerfeld | \leerfeld | \leerfeld )
\teil Neben einem Fahnenmast mit der Spitze $S(-1|1|5)$ steht ein Pfosten; sein oberes Ende $P(2|3|2)$ wirft bei parallelem Sonnenlicht den Schatten $P'(3|2|0)$ auf den Boden ($xy$-Ebene). Bestimme den Schatten $S'$ der Mastspitze. \hfill (Abitur 2020 GK) \\ S'( \leerfeld | \leerfeld | \leerfeld )
\teil Ein Rasenroboter fährt geradlinig und gleichförmig von $P(8|2|0)$ über $Q(6|3|0)$ weiter; für die Strecke von $P$ nach $Q$ braucht er 4 s. Eine Lichtschranke liegt in der senkrechten Ebene $E$: $x - 2y = -6$. Wie lange braucht der Roboter von $Q$ bis zur Lichtschranke? \hfill (Abitur 2018 GK) \\ \leerfeld[s]
\teil Ein Boot fährt geradlinig mit konstanter Geschwindigkeit von $P(1|8|0)$ nach $Q(4|6|0)$ und braucht dafür 2 Minuten; es fährt geradeaus weiter. Die Grenze eines Schutzgebiets liegt in der senkrechten Ebene $E$: $2x + y = 20$. Wie lange braucht das Boot von $Q$ bis zur Grenze? \hfill (Abitur 2018 GK) \\ \leerfeld[min]
\teil Eine Drohne fliegt geradlinig und gleichförmig von $P(2|1|4)$ über $Q(3|3|5)$ weiter; von $P$ nach $Q$ braucht sie 5 s. Eine Hauswand liegt in der Ebene $E$: $x + y = 15$. Wie lange braucht die Drohne von $Q$ bis zur Hauswand? \hfill (Abitur 2018 GK) \\ \leerfeld[s]
\end{teile}
\end{aufgabe}

% Anlauf (ohne Prüfkennung)
\begin{aufgabe}{Schnittpunkt zweier Geraden – Anlauf}
\begin{teile}
\teil Zwei Geraden im Raum werden gleichgesetzt: $(1|0|2) + r \cdot (1|2|0) = (3|4|1) + s \cdot (0|1|1)$. Was entsteht? \\ \kreuz{drei Gleichungen mit zwei Unbekannten} \\ \kreuz{zwei Gleichungen mit zwei Unbekannten} \\ \kreuz{drei Gleichungen mit drei Unbekannten}
\teil Bestimme den Schnittpunkt der Geraden $g: \vec{x} = (2|2|3) + r \cdot (1|0|2)$ und $h: \vec{x} = (3|0|3) + s \cdot (0|1|1)$. S( \leerfeld | \leerfeld | \leerfeld )
\end{teile}
\end{aufgabe}

\begin{aufgabe}{Schnittpunkt zweier Geraden – Prüfungsaufgaben}
\begin{teile}
\teil Ein Mähroboter startet auf einer ebenen, leicht geneigten Rasenfläche im Punkt $P(2|4{,}5|0{,}25)$ und fährt entlang der Geraden durch $P$ mit dem Richtungsvektor $(4|-1|0{,}5)$ auf den Rand $BC$ zu; $B(8|0|1)$, $C(8|6|1)$, eine Längeneinheit entspricht 1 m. Bestimme den Punkt $Q$, in dem die Bahn die Strecke $BC$ trifft (zur Kontrolle: $Q(8|3|1)$). \hfill (Abitur 2021 GK) \\ Q( \leerfeld | \leerfeld | \leerfeld )
\teil Gegeben sind $g: \vec{x} = (1|4|-2) + s \cdot (2|0|3)$ und die Gerade $h$ durch $A(3|0|1)$ und $B(4|1|b)$ mit einer reellen Zahl $b$. $g$ und $h$ haben einen gemeinsamen Punkt. Bestimme $b$. \hfill (Abitur 2023 GK) \\ b = \leerfeld
\teil Gegeben sind $g: \vec{x} = (0|-2|5) + s \cdot (1|0|-1)$ und die Gerade $h$ durch $A(2|1|0)$ und $B(3|2|b)$ mit einer reellen Zahl $b$. $g$ und $h$ haben einen gemeinsamen Punkt. Bestimme $b$. \hfill (Abitur 2023 GK) \\ b = \leerfeld
\teil Gegeben sind $g: \vec{x} = (3|1|0) + s \cdot (0|2|3)$ und die Gerade $h$ durch $A(1|3|2)$ und $B(2|3|b)$ mit einer reellen Zahl $b$. $g$ und $h$ haben einen gemeinsamen Punkt. Bestimme $b$. \hfill (Abitur 2023 GK) \\ b = \leerfeld
\teil Gegeben ist der Würfel $ABCDEFGH$ mit $A(0|0|0)$, $G(4|4|4)$ und $H(0|4|4)$. Für jede reelle Zahl $r$ ist die Gerade $g_r: \vec{x} = (1 - r|0|r^2) + u \cdot (2|4|0)$, $u \in \mathbb{R}$, gegeben. Für einen Wert von $r$ schneidet $g_r$ die Kante $GH$. Bestimme das Verhältnis, in dem der Schnittpunkt die Kante $GH$ teilt. \hfill (Abitur 2019 GK) \\ HQ : QG = \leerfeld : \leerfeld
\teil Gegeben ist der Würfel $ABCDEFGH$ mit $A(0|0|0)$, $G(6|6|6)$ und $H(0|6|6)$. Für jede reelle Zahl $r$ ist die Gerade $g_r: \vec{x} = (3 + r|0|r^2 - 3) + u \cdot (1|3|0)$, $u \in \mathbb{R}$, gegeben. Für einen Wert von $r$ schneidet $g_r$ die Kante $GH$. Bestimme das Verhältnis, in dem der Schnittpunkt die Kante $GH$ teilt. \hfill (Abitur 2019 GK) \\ HQ : QG = \leerfeld : \leerfeld
\teil Gegeben ist der Würfel $ABCDEFGH$ mit $A(0|0|0)$, $G(2|2|2)$ und $H(0|2|2)$. Für jede reelle Zahl $r$ ist die Gerade $g_r: \vec{x} = (1 - 0{,}5r|0|r^2 - 2) + u \cdot (1|2|0)$, $u \in \mathbb{R}$, gegeben. Für einen Wert von $r$ schneidet $g_r$ die Kante $GH$. Bestimme das Verhältnis, in dem der Schnittpunkt die Kante $GH$ teilt. \hfill (Abitur 2019 GK) \\ HQ : QG = \leerfeld : \leerfeld
\teil Gegeben ist der Würfel $ABCDEFGH$ mit $A(0|0|0)$, $G(8|8|8)$ und $H(0|8|8)$. Für jede reelle Zahl $r$ ist die Gerade $g_r: \vec{x} = (7 + 2r|0|2r^2) + u \cdot (1|4|0)$, $u \in \mathbb{R}$, gegeben. Für einen Wert von $r$ schneidet $g_r$ die Kante $GH$. Bestimme das Verhältnis, in dem der Schnittpunkt die Kante $GH$ teilt. \hfill (Abitur 2019 GK) \\ HQ : QG = \leerfeld : \leerfeld
\end{teile}
\end{aufgabe}

\begin{aufgabe}{Schattenpunkt auf einer Kante als Schnittpunkt von Lichtgerade und Kantengerade berechnen – Prüfungsaufgaben}
\begin{teile}
\teil Ein Spielturm hat ein Dach in Form einer Pyramide; an der Spitze ist eine senkrechte Stange mit dem oberen Ende $T(0|0|7)$ befestigt; der Boden ist die $xy$-Ebene, eine Längeneinheit entspricht 1 m. Sonnenlicht fällt in parallelen Geraden mit dem Richtungsvektor $(4|-1|-2)$. Der Schatten des oberen Endes fällt auf die Dachkante $EF$ mit $E(6|-3|4)$ und $F(6|3|4)$. Bestimme diesen Schattenpunkt. \hfill (Abitur 2017 GK) \\ ( \leerfeld | \leerfeld | \leerfeld )
\teil An einem Pavillon ist eine Wetterfahne mit dem oberen Ende $T(1|1|7)$ befestigt; der Boden ist die $xy$-Ebene, eine Längeneinheit entspricht 1 m. Sonnenlicht fällt in parallelen Geraden mit dem Richtungsvektor $(2|1|-3)$. Der Schatten des oberen Endes fällt auf die Dachkante $EF$ mit $E(3|-1|4)$ und $F(3|5|4)$. Bestimme diesen Schattenpunkt. \hfill (Abitur 2017 GK) \\ ( \leerfeld | \leerfeld | \leerfeld )
\teil Auf einem Gartenhaus steht ein Blitzableiter mit der Spitze $T(0|0|4{,}5)$; der Boden ist die $xy$-Ebene, eine Längeneinheit entspricht 1 m. Sonnenlicht fällt in parallelen Geraden mit dem Richtungsvektor $(-2|2|-1)$. Der Schatten des oberen Endes fällt auf die Dachkante $EF$ mit $E(-3|1|4)$ und $F(1|1|4)$. Bestimme diesen Schattenpunkt. \hfill (Abitur 2017 GK) \\ ( \leerfeld | \leerfeld | \leerfeld )
\end{teile}
\end{aufgabe}

% Anlauf (ohne Prüfkennung)
\begin{aufgabe}{Spuren, Schnittgeraden und Schnittfiguren – Anlauf}
\begin{teile}
\teil Der Spurpunkt einer Ebene auf der $y$-Achse – welche Koordinaten sind null? \\ \kreuz{$x$ und $z$} \\ \kreuz{nur $y$} \\ \kreuz{nur $x$}
\teil Bestimme die Spurpunkte der Ebene $2x + 3y + 6z = 6$ auf den Koordinatenachsen. $S_x($ \leerfeld |0|0), $S_y(0|$ \leerfeld |0), $S_z(0|0|$ \leerfeld )
\end{teile}
\end{aufgabe}

\begin{aufgabe}{Spuren, Schnittgeraden und Schnittfiguren – Prüfungsaufgaben}
\begin{teile}
\teil Bestimme die Spurpunkte der Ebene $6x + 4y + 3z = 24$ auf den Koordinatenachsen. \hfill (Abitur 2020 GK) \\ $S_x($ \leerfeld |0|0), $S_y(0|$ \leerfeld |0), $S_z(0|0|$ \leerfeld )
\teil Ein Holzkörper ist eine Pyramide über dem Quadrat $ABCD$ mit $A(0|0|0)$, $B(5|0|0)$, $C(5|5|0)$, $D(0|5|0)$ und der Spitze $E(0|5|4)$ senkrecht über $D$; eine Längeneinheit entspricht 1 cm. Die Ebene $L$: $4x + 5z = 20$ schneidet die $z$-Achse im Punkt $F$. Bestimme $F$ und zeichne die Schnittgeraden von $L$ mit der $xz$-Ebene und mit der $yz$-Ebene in das Schrägbild ein. \hfill (Abitur 2021 GK)

\begin{ksys3}[xyz,x1min=0,x1max=6,x2min=0,x2max=6,x3min=0,x3max=5] \rpyramide{0,0,0}{5}{5}{0,5,4} \end{ksys3}
\teil Ein Holzkörper ist eine Pyramide über dem Quadrat $ABCD$ mit $A(0|0|0)$, $B(7|0|0)$, $C(7|7|0)$, $D(0|7|0)$ und der Spitze $E(0|7|3)$ senkrecht über $D$; eine Längeneinheit entspricht 1 cm. Die Ebene $L$: $3x + 7z = 21$ schneidet die $z$-Achse im Punkt $F$. Bestimme $F$ und zeichne die Schnittgeraden von $L$ mit der $xz$-Ebene und mit der $yz$-Ebene in das Schrägbild ein. \hfill (Abitur 2021 GK)

\begin{ksys3}[xyz,x1min=0,x1max=8,x2min=0,x2max=8,x3min=0,x3max=4] \rpyramide{0,0,0}{7}{7}{0,7,3} \end{ksys3}
\teil Die Abbildung zeigt den Quader mit $A(7|0|0)$, $B(7|7|0)$, $C(0|7|0)$ und $F(7|7|2)$, die Gerade $h$ durch $B$ und $F$ sowie die Schnittfigur des Quaders mit einer Ebene durch $A$, $C$ und einen Punkt $P$ von $h$. Beschreibe, wie man mithilfe der Abbildung ermitteln kann, dass $P$ die $z$-Koordinate 4 hat. \hfill (Abitur 2018 GK)

\begin{ksys3}[xyz,x1min=0,x1max=8,x2min=0,x2max=8,x3min=0,x3max=5] \rquader{0,0,0}{7}{7}{2} \rgerade[0:1]{7,0,0}{0,3.5,2}{} \rgerade[0:1]{7,3.5,2}{-3.5,3.5,0}{} \rgerade[0:1]{3.5,7,2}{-3.5,0,-2}{} \rgerade[0:1]{0,7,0}{7,-7,0}{} \rgerade[0:5]{7,7,0}{0,0,1}{h} \end{ksys3}
\teil Die Abbildung zeigt den Quader mit $A(5|0|0)$, $B(5|5|0)$, $C(0|5|0)$ und $F(5|5|2)$, die Gerade $h$ durch $B$ und $F$ sowie die Schnittfigur des Quaders mit einer Ebene durch $A$, $C$ und einen Punkt $P$ von $h$. Beschreibe, wie man mithilfe der Abbildung ermitteln kann, dass $P$ die $z$-Koordinate 4 hat. \hfill (Abitur 2018 GK)

\begin{ksys3}[xyz,x1min=0,x1max=6,x2min=0,x2max=6,x3min=0,x3max=5] \rquader{0,0,0}{5}{5}{2} \rgerade[0:1]{5,0,0}{0,2.5,2}{} \rgerade[0:1]{5,2.5,2}{-2.5,2.5,0}{} \rgerade[0:1]{2.5,5,2}{-2.5,0,-2}{} \rgerade[0:1]{0,5,0}{5,-5,0}{} \rgerade[0:5]{5,5,0}{0,0,1}{h} \end{ksys3}
\end{teile}
\end{aufgabe}

\begin{aufgabe}{Spuren, Schnittgeraden und Schnittfiguren – Prüfungsaufgaben – weiter}
\begin{teile}
\teil Gegeben ist das Quadrat $ABCD$ mit $A(0|0|0)$, $B(8|0|0)$, $C(8|8|0)$ und $D(0|8|0)$ sowie die Ebenen $E_1$: $x = 4$, $E_2$: $x + y = 8$ und $E_3$: $y = 2$. Entscheide für jede der drei Ebenen, ob sie das Quadrat in zwei Figuren mit gleichem Flächeninhalt teilt, und begründe. \hfill (Abitur 2021 GK)
\teil Gegeben ist das Quadrat $ABCD$ mit $A(0|0|0)$, $B(5|0|0)$, $C(5|5|0)$ und $D(0|5|0)$ sowie die Ebenen $E_1$: $z = 1$, $E_2$: $x - y = 0$ und $E_3$: $x = 2$. Entscheide für jede der drei Ebenen, ob sie das Quadrat in zwei Figuren mit gleichem Flächeninhalt teilt, und begründe. \hfill (Abitur 2021 GK)
\teil Gegeben ist das Quadrat $ABCD$ mit $A(1|1|0)$, $B(7|1|0)$, $C(7|7|0)$ und $D(1|7|0)$ sowie die Ebenen $E_1$: $y = 4$, $E_2$: $x + y = 8$ und $E_3$: $x = 5$. Entscheide für jede der drei Ebenen, ob sie das Quadrat in zwei Figuren mit gleichem Flächeninhalt teilt, und begründe. \hfill (Abitur 2021 GK)
\teil Gegeben sind die Ebenen $E_1: \vec{x} = (5|0|0) + r \cdot (-1|1|0) + s \cdot (-1|0|1)$ und $E_2$: $2x + 2y + z = 10$. Begründe, dass $E_1$ und $E_2$ nicht parallel sind, und bestimme eine Gleichung ihrer Schnittgeraden. \hfill (Abitur 2020 GK) \\ x = \leerfeld + r · \leerfeld
\teil Gegeben sind die Ebenen $E_1: \vec{x} = (0|3|0) + r \cdot (2|-3|0) + s \cdot (0|-3|1)$ und $E_2$: $2x + y + 4z = 4$. Begründe, dass $E_1$ und $E_2$ nicht parallel sind, und bestimme eine Gleichung ihrer Schnittgeraden. \hfill (Abitur 2020 GK) \\ x = \leerfeld + r · \leerfeld
\end{teile}
\end{aufgabe}

\clearpage
% Abstände – bank/abstaende/ (zusammenbau v0.4, Heft)
\einheitenkopf[t7]{Abstände}
\setcounter{aufgabe}{64}

% Anlauf (ohne Prüfkennung)
\begin{aufgabe}{Abstand zweier Punkte – Anlauf}
\begin{teile}
\teil Abstand wovon zu was? Wie weit ist die Mastspitze S(3 | 2 | 7) vom Punkt T(0 | 6 | 7) entfernt? Kreuze das Verfahren an. \\ \kreuz{Punkt–Punkt: Betrag des Verbindungsvektors} \\ \kreuz{Punkt–Ebene: Hessesche Normalform} \\ \kreuz{Punkt–Gerade: Lotfußpunkt}
\teil A(2 | 1 | 3), B(4 | 3 | 4) – wie weit voneinander entfernt? d = \leerfeld
\teil Ein Segelflugzeug gleitet geradlinig in 40 s von A(200 | 300 | 1\,000) nach B(800 | 1\,500 | 600); 1 LE = 1 m. Wie schnell ist es in km/h? v = \leerfeld km/h
\end{teile}
\end{aufgabe}

\begin{aufgabe}{Abstand zweier Punkte – Prüfungsaufgaben}
\begin{teile}
\teil Eine Drohne fliegt geradlinig in 4 s von P(5 | 20 | 30) nach Q(65 | 100 | 30); 1 LE = 1 m. Wie schnell ist sie in km/h? \hfill (Abitur 2018 GK) \\ v = \leerfeld km/h
\teil Ein Klettergerüst hat die Kanten AB mit A(0 | 1 | 4), B(4 | 1 | 4) und CD mit C(2 | 7 | 1), D(6 | 7 | 1); 1 LE = 1 m. Ein Seil zwischen den Mittelpunkten der beiden Kanten ist 30 \% länger als deren Abstand. Wie lang ist das Seil? \hfill (Abitur 2018 GK) \\ \leerfeld[m]
\teil Eine rechteckige Terrasse hat die Ecken A(2 | 1 | 0), B(9 | 1 | 0), C(9 | 7 | 0) und D(2 | 7 | 0); 1 LE = 1 m. Ein Gartenschlauch ist am Wasserhahn P(3 | 2 | 2) befestigt und 9 m lang. Erreicht er jede Stelle der Terrasse? \hfill (Abitur 2023 GK) \\ \janein
\teil Das Dreieck ABC mit A(0 | 0 | 0), B(2 | 1 | 2) und C(t | −2t | 0), t positiv, ist rechtwinklig in A und gleichschenklig. Welcher Wert von t? \hfill (Abitur 2026 GK) \\ t = \leerfeld
\teil Gegeben sind P(1 | 3 | −2) und Q(9 | y | 13) mit unbekannter Koordinate y. Weise nach, dass P und Q mindestens den Abstand 17 haben. \hfill (Abitur 2019 GK)
\teil Gegeben sind R(−1 | 4 | 2) und S(4 | −8 | z) mit unbekannter Koordinate z. Weise nach, dass R und S mindestens den Abstand 13 haben. \hfill (Abitur 2019 GK)
\end{teile}
\end{aufgabe}

% Anlauf (ohne Prüfkennung)
\begin{aufgabe}{Abstand Punkt–Ebene – Anlauf}
\begin{teile}
\teil Hin oder zurück? Wie weit ist P(2 | −1 | 3) von der Ebene E: $2x - 2y + z = 1$ entfernt? Kreuze an, rechne nicht. \\ \kreuz{Abstand berechnen} \\ \kreuz{Punkt zum vorgegebenen Abstand bestimmen}
\teil Abstand von P(4 | 0 | 3) zur Ebene E: $2x - y + 2z = 5$? d = \leerfeld
\teil Lotgerade durch P(−1 | 0 | 5) zur Ebene E: $3x - 4z = 1$? $\vec x = $ \leerfeld
\end{teile}
\end{aufgabe}

\begin{aufgabe}{Abstand Punkt–Ebene – Prüfungsaufgaben}
\begin{teile}
\teil Die Ebene E: $3y + 4z = 3$ enthält den Punkt (2 | 1 | 0). Gesucht ist ein Punkt P, dessen Spiegelpunkt an E von P den Abstand 14 hat. Gib einen solchen Punkt an. \hfill (Abitur 2025 GK) \\ P( \leerfeld | \leerfeld | \leerfeld )
\end{teile}
\end{aufgabe}

% Anlauf (ohne Prüfkennung)
\begin{aufgabe}{Lotfußpunkt – Anlauf}
\begin{teile}
\teil Was sagt die Gleichung $\overrightarrow{PF} \circ \vec r = 0$ in einem Lösungsweg? Kreuze an. \\ \kreuz{Punkt auf der Geraden} \\ \kreuz{Lotbedingung (Skalarprodukt null)} \\ \kreuz{Abstandsgleichheit}
\teil Lotfußpunkt F von P(4 | 2 | 3) auf g: $\vec x = (2 | 1 | 0) + t \cdot (2 | -1 | 2)$ – und der Abstand von P zu g? F( \leerfeld | \leerfeld | \leerfeld ), d ≈ \leerfeld
\teil Die Gleichungen I $\overrightarrow{OX} = \overrightarrow{OM} + s \cdot \vec u$ und II $\overrightarrow{PX} \circ \vec u = 0$ liefern gemeinsam einen Wert von s. Welche Bedeutung haben X und die Länge $|\overrightarrow{PX}|$?
\end{teile}
\end{aufgabe}

\begin{aufgabe}{Lotfußpunkt – Prüfungsaufgaben}
\begin{teile}
\teil An einem Zeltgestänge liefern die Gleichungen I $\overrightarrow{OQ} = \overrightarrow{OA} + r \cdot \overrightarrow{AB}$ und II $\overrightarrow{EQ} \circ \overrightarrow{AB} = 0$ gemeinsam einen Wert von r. Welche Bedeutung hat der Punkt Q für diesen Wert? \hfill (Abitur 2026 GK)
\teil Bei einem Carport soll von der Mitte M des Balkens AB eine möglichst kurze Strebe zum Balken CD führen; sie trifft CD im Punkt R. Beschreibe, wie man die Koordinaten von R ermitteln kann. \hfill (Abitur 2022 GK)
\teil Die Kante AB eines Körpers liegt in der xy-Ebene und in der Ebene L: $2x + y + z = 8$. Zum Punkt Q(0 | 3 | 0) ist vorgelegt: I g: $\vec x = (0 | 3 | 0) + t \cdot (2 | 1 | 0)$; II $2 \cdot 2t + 3 + t = 8 \Rightarrow t = 1 \Rightarrow R(2 | 4 | 0)$; III $\sqrt{(0 - 2)^2 + (3 - 4)^2} \approx 2{,}24$. Erläutere die Schritte und gib die geometrische Bedeutung des Ergebnisses an. \hfill (Abitur 2023 GK)
\teil Pyramide über dem Quadrat ABCD mit A(1 | 0 | 0), B(5 | 0 | 0), C(5 | 4 | 0), D(1 | 4 | 0) und der Spitze E(1 | 4 | 2); B, D, E liegen in einer Symmetrieebene. Vorgelegt ist ein Rechenweg für die kürzeste Linie von A über die Kante BE nach C: $P(5 - 4t | 4t | 2t)$; $\overrightarrow{PA} \circ \overrightarrow{PB} = 0 \Leftrightarrow t = \tfrac{4}{9}$; $2 \cdot |\overrightarrow{PA}| \approx 5{,}96$. Erläutere das Vorgehen. \hfill (Abitur 2021 GK)
\teil Ein Mähroboter mit kreisförmiger Unterseite (Radius 0,3 m) fährt mit seinem Mittelpunkt auf der Geraden g mit Richtungsvektor (3 | 4 | 0) auf den Rand k der Rasenfläche zu; g trifft k im Punkt Q(6 | 8 | 0) unter dem Winkel 30°; 1 LE = 1 m. Der Roboter wendet, sobald seine Unterseite k berührt. Wo ist dann sein Mittelpunkt S? Nutze eine Skizze. \hfill (Abitur 2021 GK) \\ S( \leerfeld | \leerfeld | \leerfeld )
\end{teile}
\end{aufgabe}

% Anlauf (ohne Prüfkennung)
\begin{aufgabe}{Abstand als Argument – Anlauf}
\begin{teile}
\teil Senkrecht oder schräg? Die Strecke verbindet P(1 | 2 | 4) mit dem Punkt Q(1 | 2 | 0) der xy-Ebene. Kreuze an, rechne nicht. \\ \kreuz{senkrecht: die Strecke ist der Abstand} \\ \kreuz{schräg: die Strecke ist länger als der Abstand}
\teil Die zwei Scheiben eines Doppelfensters liegen in den parallelen Ebenen $E_1$: $2x - y + 2z = 1$ und $E_2$: $2x - y + 2z = 7$; P(0 | −1 | 0) liegt in $E_1$, Q(1 | 1 | 3) in $E_2$. Ist PQ kleiner, größer oder genauso groß wie der Abstand der Scheiben? Begründe.
\teil Die Abbildung zeigt eine Rampe im Querschnitt als Strecke GJ auf der Geraden t und eine Lampe F. F hat von t etwa den Abstand 3,5. Begründe mit einer Zeichnung, dass kein Punkt der Rampe von F den Abstand 3,5 hat.

\begin{ksys}[xmin=0,xmax=8,ymin=0,ymax=8]
\gerade{-1}{8}{t}
\punkt{4}{4}{G}
\punkt{7}{1}{J}
\punkt{0}{3}{F}
\end{ksys}
\end{teile}
\end{aufgabe}

\begin{aufgabe}{Abstand als Argument – Prüfungsaufgaben}
\begin{teile}
\teil B(1 | 0 | 4) liegt in der Ebene $L_1$: $x + 2y + 2z = 9$, E(1 | 0 | 1) in der parallelen Ebene $L_2$: $x + 2y + 2z = 3$. Ist der Abstand von B und E kleiner, größer oder genauso groß wie der Abstand der Ebenen? Begründe. \hfill (Abitur 2023 GK)
\teil Ein Dachfirstpunkt F ist 3 m vom Punkt G einer schrägen Glasfläche entfernt; die Strecke FG steht nicht senkrecht auf der Glasfläche. Begründe, dass der Abstand von F zur Ebene der Glasfläche kleiner als 3 m ist. \hfill (Abitur 2023 GK)
\teil Eine Seilrutsche führt vom Startpunkt $S_1$ zum tiefer gelegenen Zielpunkt $S_2$. Auf der Karte erscheint als Länge der Anlage der horizontale Abstand von $S_1$ und $S_2$. Begründe ohne Rechnung, dass diese Länge kürzer ist als die Seilstrecke $S_1S_2$. \hfill (Abitur 2021 GK)
\teil Gegeben sind P(2 | 1 | 0), Q(8 | 1 | 0) und R(2 | 5 | 0). Gib zwei Punkte an, die von P, Q und R gleich weit entfernt sind. \hfill (Abitur 2019 GK) \\ ( \leerfeld | \leerfeld | \leerfeld ), ( \leerfeld | \leerfeld | \leerfeld )
\teil Ein Zeltdach hat die Spitze S(2 | 2 | 9); F(3 | 2 | 6) und G(2 | 3 | 6) liegen auf zwei Dachkanten. Die Geraden SF und SG schneiden die xy-Ebene in $Q_1$ und $Q_2$. In welchem Verhältnis steht der Abstand von $Q_1$ und $Q_2$ zum Abstand von F und G? Begründe, ohne $Q_1$ und $Q_2$ zu berechnen. \hfill (Abitur 2022 GK)
\teil Eine Pyramide hat die Spitze S(1 | 1 | 8); F(2 | 1 | 6) und G(1 | 2 | 6) liegen auf zwei Seitenkanten. Die Geraden SF und SG schneiden die xy-Ebene in $Q_1$ und $Q_2$. In welchem Verhältnis steht der Abstand von $Q_1$ und $Q_2$ zum Abstand von F und G? Begründe, ohne $Q_1$ und $Q_2$ zu berechnen. \hfill (Abitur 2022 GK)
\end{teile}
\end{aufgabe}

\clearpage
% Skalarprodukt und Winkel – bank/skalarprodukt-und-winkel/ (zusammenbau v0.4, Heft)
\einheitenkopf[t8]{Skalarprodukt und Winkel}
\setcounter{aufgabe}{72}

% Anlauf (ohne Prüfkennung)
\begin{aufgabe}{Das Skalarprodukt als Objekt – Anlauf}
\begin{teile}
\teil $\vec a$, $\vec b$ und $\vec c$ sind Vektoren mit je drei Koordinaten. Was ergibt $\vec a \circ \vec b$? Kreuze an. \\ \kreuz{Vektor} \\ \kreuz{Zahl} \\ \kreuz{nicht definiert}
\teil $\vec a = (2 \mid 1 \mid 3)$, $\vec b = (1 \mid 4 \mid -1)$. Berechne $\vec a \circ \vec b$ und gib an, ob der Winkel zwischen $\vec a$ und $\vec b$ spitz, recht oder stumpf ist. $\vec a \circ \vec b =$ \leerfeld \quad Winkel: \leerfeld
\teil $\vec a$, $\vec b$ und $\vec c$ sind Vektoren mit je drei Koordinaten. Prüfe die Rechenarten von innen nach außen: Was ergibt $(\vec a + \vec b) \circ \vec c$? \\ \kreuz{Vektor} \\ \kreuz{Zahl} \\ \kreuz{nicht definiert}
\end{teile}
\end{aufgabe}

\begin{aufgabe}{Das Skalarprodukt als Objekt – Prüfungsaufgaben}
\begin{teile}
\teil Ein Carsharing-Anbieter verteilt seine $600$ Fahrzeuge auf drei Stadtteile. Der Verteilungsvektor $(x \mid y \mid z)$ gibt die Anzahlen an, es gilt also $x + y + z = 600$. Beurteile beide Aussagen: (1) Es gibt zwei verschiedene Verteilungen, deren Vektoren den Winkel $0^\circ$ einschließen. (2) Es gibt zwei Verteilungen, deren Vektoren den Winkel $90^\circ$ einschließen. \hfill (Abitur 2024 GK)
\teil Ein Quader hat eine quadratische Grundfläche mit der Seitenlänge s und die Höhe h. An einer Ecke der Grundfläche sind $\vec u$ und $\vec v$ die Kantenvektoren der Grundfläche und $\vec w$ der Kantenvektor der Höhe. Begründe ohne Koordinaten, dass $\vec u \circ (\vec u + \vec v + 2\vec w)$ nur von s abhängt. \hfill (Abitur 2021 GK)
\teil Ein Quader hat eine quadratische Grundfläche mit der Seitenlänge s und die Höhe h. An einer Ecke der Grundfläche sind $\vec u$ und $\vec v$ die Kantenvektoren der Grundfläche und $\vec w$ der Kantenvektor der Höhe. $\vec e = \vec u + \vec v$ ist der Vektor der Diagonalen der Grundfläche. Begründe ohne Koordinaten, dass $\vec e \circ (\vec e + 3\vec w)$ nur von s abhängt. \hfill (Abitur 2021 GK)
\end{teile}
\end{aufgabe}

% Anlauf (ohne Prüfkennung)
\begin{aufgabe}{Winkel zwischen Vektoren, Kanten und Geraden – Anlauf}
\begin{teile}
\teil Gesucht ist der Innenwinkel des Dreiecks ABC bei A. Welche zwei Vektoren setzt du an? \\ \kreuz{$\overrightarrow{AB}$ und $\overrightarrow{AC}$} \\ \kreuz{$\overrightarrow{BA}$ und $\overrightarrow{AC}$} \\ \kreuz{$\overrightarrow{AB}$ und $\overrightarrow{BC}$}
\teil Gegeben sind $P(1 \mid 2 \mid 0)$, $Q(3 \mid 3 \mid 2)$ und $R(4 \mid 0 \mid 1)$. Berechne den Winkel zwischen den Kanten PQ und PR. $\varphi \approx$ \leerfeld[$^\circ$]
\teil Gegeben sind $A(2 \mid 1 \mid 0)$, $B(4 \mid 3 \mid 1)$ und $C(5 \mid 1 \mid 3)$. Weise nach, dass das Dreieck ABC rechtwinklig ist, und gib die Ecke mit dem rechten Winkel an.
\end{teile}
\end{aufgabe}

\begin{aufgabe}{Winkel zwischen Vektoren, Kanten und Geraden – Prüfungsaufgaben}
\begin{teile}
\teil Ein Zelt hat die Bodenecken $K(5 \mid 2 \mid 0)$, $L(1 \mid 3 \mid 0)$ und $S(3 \mid 1 \mid 6)$. S ist die Zeltspitze. Berechne die Größe des Winkels zwischen den Kanten KL und KS. \hfill (Abitur 2026 GK) \\ $\varphi \approx$ \leerfeld[$^\circ$]
\teil Ein Würfel mit der Kantenlänge 7 hat eine Ecke im Ursprung, seine Kanten dort liegen auf den positiven Achsen. Auf seinen Kanten liegen $I(7 \mid 0 \mid 1)$, $J(3 \mid 7 \mid 0)$, $K(0 \mid 7 \mid 3)$ und $L(1 \mid 0 \mid 7)$. Das Viereck IJKL ist ein Trapez. Berechne die Größe des Innenwinkels bei I. \hfill (Abitur 2019 GK) \\ $\varphi \approx$ \leerfeld[$^\circ$]
\teil Eine Pyramide ABCDS hat ihre Grundfläche in der xy-Ebene; bekannt sind $C(5 \mid 4 \mid 0)$, $D(1 \mid 8 \mid 0)$ und $S(1 \mid 0 \mid 3)$. Weise nach, dass die Seitenfläche CDS ein rechtwinkliges Dreieck ist. \hfill (Abitur 2024 GK)
\teil Ein Saugroboter fährt vom Punkt $P(1 \mid 5 \mid 0)$ aus geradlinig in Richtung $(4 \mid -1 \mid 0)$ auf eine Wand zu, die am Boden entlang der Strecke von $B(9 \mid 0 \mid 0)$ nach $C(6 \mid 5 \mid 0)$ verläuft; 1 LE entspricht 1 m. Weise nach, dass sich der Roboter der Wand unter einem Winkel von $45^\circ$ nähert. \hfill (Abitur 2021 GK)
\end{teile}
\end{aufgabe}

% Anlauf (ohne Prüfkennung)
\begin{aufgabe}{Neigungswinkel von Ebenen – Anlauf}
\begin{teile}
\teil Der Winkel zwischen zwei Ebenen wird über ihre Normalenvektoren berechnet. Welche Funktion steht in der Formel? \\ \kreuz{Kosinus} \\ \kreuz{Sinus}
\teil Berechne den Winkel zwischen der Ebene $E\colon x + 2y + 2z = 8$ und der xy-Ebene. $\varphi \approx$ \leerfeld[$^\circ$]
\teil Zwei Seitenflächen eines Körpers liegen in den Ebenen $E\colon 2x + y + 2z = 10$ und $F\colon x - 2y + 2z = 3$. Berechne den Schnittwinkel der beiden Ebenen. $\varphi \approx$ \leerfeld[$^\circ$]
\end{teile}
\end{aufgabe}

\begin{aufgabe}{Neigungswinkel von Ebenen – Prüfungsaufgaben}
\begin{teile}
\teil Ein Hang liegt in der Ebene $E\colon 2x + 3y + z = 9$, die xy-Ebene ist die Horizontale. Berechne die Größe des Winkels zwischen Hang und Horizontale. \hfill (Abitur 2025 GK) \\ $\varphi \approx$ \leerfeld[$^\circ$]
\teil Die Pyramide ABCDS hat die Ecken $A(0 \mid 0 \mid 0)$, $B(7 \mid 0 \mid 0)$, $C(7 \mid 4 \mid 0)$, $D(0 \mid 7 \mid 0)$ und die Spitze $S(0 \mid 0 \mid 6)$. Die Seitenfläche BCS liegt in der Ebene $E\colon 6x + 7z = 42$. Berechne den Winkel zwischen E und der Ebene F, in der die Seitenfläche ADS liegt. \hfill (Abitur 2024 GK) \\ $\varphi \approx$ \leerfeld[$^\circ$]
\teil Eine Kletterwand liegt in der Ebene $L\colon 6x + 3y + 2z = 24$, der Untergrund ist die xy-Ebene; 1 LE entspricht 1 m. Berechne die Größe des Winkels zwischen Kletterwand und Untergrund. \hfill (Abitur 2018 GK) \\ $\varphi \approx$ \leerfeld[$^\circ$]
\teil Ein Hausdach liegt in der Ebene $3y - 4z = 0$, die Traufe (untere Dachkante) liegt auf der x-Achse. Von der Traufe aus steigt das Dach über dem Haus (Bereich $y > 0$) an, die Hauswand verläuft von der Traufe aus senkrecht nach unten. Berechne den Neigungswinkel des Dachs gegen die Horizontale und den Winkel im Inneren des Hauses zwischen Dach und Wand. \hfill (Abitur 2024 GK) \\ $\varphi \approx$ \leerfeld[$^\circ$] \quad Innenwinkel \leerfeld[$^\circ$]
\teil Ein Sonnensegel liegt in der Ebene $3y + 8z = 24$, der Untergrund ist die xy-Ebene; 1 LE entspricht 1 m. Damit Regenwasser abläuft, muss die Neigung des Segels mindestens $35\,\%$ betragen. Prüfe, ob diese Bedingung erfüllt ist. \hfill (Abitur 2020 GK)
\teil Eine Kellerklappe ist an der x-Achse drehbar befestigt und steht geöffnet in der Ebene $2y - 3z = 0$; der Boden ist die xy-Ebene. Beim Schließen wird die Klappe gedreht, bis der keilförmige Raum unter ihr kein Volumen mehr hat. Berechne den Drehwinkel und erläutere den Ansatz. \hfill (Abitur 2026 GK)
\teil Eine Zugbrücke ist an der x-Achse drehbar gelagert; hochgezogen liegt sie in der Ebene $3y - z = 0$, die Fahrbahn davor ist die xy-Ebene. Die Brücke wird herabgelassen, bis sie waagerecht aufliegt. Berechne den Drehwinkel und erläutere den Ansatz. \hfill (Abitur 2026 GK)
\end{teile}
\end{aufgabe}

% Anlauf (ohne Prüfkennung)
\begin{aufgabe}{Winkel gegen Ebenen und Bogenmaße – Anlauf}
\begin{teile}
\teil Der Winkel zwischen einer Treppenkante und dem Boden wird über den Richtungsvektor der Kante und einen Normalenvektor des Bodens berechnet. Welche Funktion steht in der Formel? \\ \kreuz{Kosinus} \\ \kreuz{Sinus}
\teil Berechne den Winkel, unter dem die Gerade $g\colon \vec x = (1 \mid 0 \mid 2) + t \cdot (2 \mid 1 \mid 2)$, $t \in \mathbb{R}$, die xy-Ebene schneidet. $\varphi \approx$ \leerfeld[$^\circ$]
\teil Ein Pyramidendach hat die Traufecke $P(3 \mid 3 \mid 2)$ und die Spitze $S(0 \mid 0 \mid 5)$; die Grundfläche des Hauses liegt in der xy-Ebene. Berechne den Neigungswinkel der Dachkante PS gegen die Grundfläche. $\varphi \approx$ \leerfeld[$^\circ$]
\end{teile}
\end{aufgabe}

\begin{aufgabe}{Winkel gegen Ebenen und Bogenmaße – Prüfungsaufgaben}
\begin{teile}
\teil Sonnenlicht fällt in parallelen Geraden mit dem Richtungsvektor $(2 \mid 1 \mid -3)$ auf den ebenen Untergrund, die xy-Ebene. Berechne den Winkel, unter dem das Licht auf den Untergrund trifft. \hfill (Abitur 2023 GK) \\ $\varphi \approx$ \leerfeld[$^\circ$]
\end{teile}
\end{aufgabe}

\clearpage
% Orthogonalität – bank/orthogonalitaet/ (zusammenbau v0.4, Heft)
\einheitenkopf[t9]{Orthogonalität}
\setcounter{aufgabe}{80}

% Anlauf (ohne Prüfkennung)
\begin{aufgabe}{Rechte Winkel nachweisen – Anlauf}
\begin{teile}
\teil Das Dreieck $ABC$ soll bei $B$ einen rechten Winkel haben. Welche Vektoren prüfst du? \\ \kreuz{$\overrightarrow{BA}$ und $\overrightarrow{BC}$} \\ \kreuz{$\overrightarrow{AB}$ und $\overrightarrow{AC}$} \\ \kreuz{$\overrightarrow{CA}$ und $\overrightarrow{CB}$}
\teil Gegeben sind $A(2 | 1 | 3)$, $B(5 | 3 | 3)$, $C(0 | 4 | 3)$. Weise nach, dass das Dreieck bei $A$ einen rechten Winkel hat.
\teil Gegeben sind $P(4 | 0 | 0)$, $Q(0 | 3 | 5)$ und der Ursprung $O$. Begründe ohne Skalarprodukt, dass das Dreieck $OPQ$ bei $O$ rechtwinklig ist.
\end{teile}
\end{aufgabe}

\begin{aufgabe}{Rechte Winkel nachweisen – Prüfungsaufgaben}
\begin{teile}
\teil Das Dreieck $OAB$ mit $O(0 | 0 | 0)$, $A(5 | 12 | 0)$ und $B(0 | 0 | 13)$ ist gleichschenklig. Weise nach, dass es bei $O$ rechtwinklig ist und dass beide Katheten die Länge 13 haben \hfill (Abitur 2026 GK).
\teil Ein Holzkörper hat die Eckpunkte $A(0 | 0 | 0)$, $B(9 | 0 | 0)$, $C(9 | 6 | 0)$, $D(0 | 6 | 0)$ und $E(0 | 6 | 4)$; er ist eine Pyramide über dem Rechteck $ABCD$ mit der Spitze $E$ senkrecht über $D$ (1 LE = 1 cm). Weise nach, dass das Dreieck $BCE$ rechtwinklig ist, und berechne den Inhalt der gesamten Oberfläche \hfill (Abitur 2021 GK).
\teil Ein Holzkörper hat die Eckpunkte $A(0 | 0 | 0)$, $B(7 | 0 | 0)$, $C(7 | 7 | 0)$, $D(0 | 7 | 0)$ und $E(7 | 0 | 5)$; er ist eine Pyramide über dem Quadrat $ABCD$ mit der Spitze $E$ senkrecht über $B$ (1 LE = 1 cm). Weise nach, dass das Dreieck $CDE$ bei $C$ rechtwinklig ist, gib die Kathetenlängen an und berechne den Inhalt der Oberfläche des Körpers \hfill (Abitur 2021 GK).
\teil Gegeben sind $A(0 | 0 | 0)$, $B(4 | 3 | 5)$ und $C_t(3t | -4t | 0)$ mit $t > 0$. Weise nach, dass das Dreieck $ABC_t$ für jedes $t$ bei $A$ einen rechten Winkel hat \hfill (Abitur 2026 GK).
\teil Gegeben sind $A(1 | 6 | 2)$, $B(4 | 6 | 0)$ und $C_t(4 + 2t | 6 | 3t)$ mit $t \ne 0$. Weise nach, dass jedes Dreieck $ABC_t$ bei $B$ einen rechten Winkel hat \hfill (Abitur 2018 GK).
\teil Ein gerades Prisma $ABCDEF$ hat die Grundfläche $ABC$ mit $A(0 | -6 | 0)$, $B(\sqrt{7} | 0 | 0)$ und $C(0 | 6 | 0)$. Weise nach, dass das Dreieck $ABC$ bei $B$ nicht rechtwinklig ist \hfill (Abitur 2019 GK).
\teil Ein gerades Prisma hat die Grundfläche $ABC$ mit $A(0 | -2 | 0)$, $B(\sqrt{5} | 0 | 0)$ und $C(0 | 3 | 0)$. Weise nach, dass das Dreieck $ABC$ bei $B$ nicht rechtwinklig ist \hfill (Abitur 2019 GK).
\end{teile}
\end{aufgabe}

% Anlauf (ohne Prüfkennung)
\begin{aufgabe}{Rechte Winkel rückwärts – Anlauf}
\begin{teile}
\teil Aufgabe: „Bestimme $a$ so, dass das Dreieck $OAB$ bei $B$ rechtwinklig ist.“ Was ist zu tun? \\ \kreuz{null zeigen – für jeden Wert des Parameters} \\ \kreuz{null setzen – ein Wert ist gesucht, Gleichung lösen} \\ \kreuz{null zeigen – ohne Parameter, einsetzen und ausrechnen}
\teil Gegeben sind $A(1 | 2 | 3)$, $B(4 | 3 | 3)$ und $C(2 | t | 5)$. Bestimme $t$ so, dass das Dreieck $ABC$ bei $A$ rechtwinklig ist.
\teil Ein Quader hat die Eckpunkte $O(0 | 0 | 0)$, $A(12 | 0 | 0)$, $B(12 | 9 | 0)$, $C(0 | 9 | 0)$ und $E(12 | 0 | h)$ mit $h > 0$; $G$ liegt über $C$. Die Raumdiagonalen $AG$ und $CE$ stehen senkrecht zueinander. Bestimme $h$ sowie Volumen und Oberflächeninhalt des Quaders.
\end{teile}
\end{aufgabe}

\begin{aufgabe}{Rechte Winkel rückwärts – Prüfungsaufgaben}
\begin{teile}
\teil Gegeben sind der Ursprung $O$, $A(4 | 0 | a)$ und $B(2 | 2 | 3)$. Bestimme $a$ so, dass das Dreieck $OAB$ bei $B$ rechtwinklig ist \hfill (Abitur 2021 GK).
\teil Der Ursprung $O$, $P(4 | 0 | 12)$ und $Q_t(2 | t | 8)$ bilden ein Dreieck. Bestimme alle Werte von $t$, für die das Dreieck bei $Q_t$ einen rechten Winkel hat \hfill (Abitur 2024 GK).
\teil Gegeben sind $M(5 | 1 | 2)$, $F(0 | 0 | 2)$ und $S_t(t | 0 | 0)$. Bestimme alle $t$, für die das Dreieck $MFS_t$ bei $S_t$ einen rechten Winkel hat \hfill (Abitur 2024 GK).
\teil Gegeben sind $A(1 | 2 | 0)$ und $B(4 | 4 | 2)$. Für den Punkt $C$ gilt: Seine $z$-Koordinate ist $5$, die Geraden $AB$ und $AC$ verlaufen senkrecht zueinander, und $C$ liegt in der Ebene $2x + y = 0$. Bestimme die Koordinaten von $C$ \hfill (Abitur 2026 GK).
\teil Gegeben sind $A(0 | 1 | 2)$ und $B(1 | 4 | 1)$. Für den Punkt $C$ gilt: Seine $z$-Koordinate ist $-1$, die Geraden $AB$ und $AC$ verlaufen senkrecht zueinander, und $C$ liegt in der Ebene $x + y = 2$. Bestimme die Koordinaten von $C$ \hfill (Abitur 2026 GK).
\end{teile}
\end{aufgabe}

% Anlauf (ohne Prüfkennung)
\begin{aufgabe}{Senkrecht zu Geraden und Ebenen – Anlauf}
\begin{teile}
\teil Aufgabe: „Zeige, dass die Gerade $h$ senkrecht auf der Ebene $E$ steht.“ Welche Prüfregel? \\ \kreuz{zwei Richtungen – Skalarprodukt null} \\ \kreuz{Gerade gegen Ebene – Richtungsvektor kollinear zum Normalenvektor} \\ \kreuz{Ebene gegen Ebene – Skalarprodukt der Normalenvektoren null}
\teil Gegeben sind $E\colon x + y + z = 1$ und $F\colon 2x - ty + 3z = 0$. Bestimme $t$ so, dass $E$ und $F$ senkrecht zueinander stehen.
\end{teile}
\end{aufgabe}

\begin{aufgabe}{Senkrecht zu Geraden und Ebenen – Prüfungsaufgaben}
\begin{teile}
\teil Ein Prisma $ABCDEF$ hat die Grundfläche $ABC$ in der Ebene $L\colon 2x - y + 2z = 6$ mit $A(3 | 0 | 0)$, $B(1 | 2 | 3)$, $C(0 | -2 | 2)$; die Seitenkante $BE$ führt zu $E(5 | 0 | 7)$. Begründe, dass das Prisma gerade ist \hfill (Abitur 2020 GK).
\teil Gegeben sind $E\colon \vec x = (2 | 0 | 1) + r \cdot (2 | 1 | 2) + s \cdot (1 | 0 | -1)$ und der Vektor $\vec n = (1 | -4 | 1)$. Weise nach, dass $\vec n$ senkrecht zu $E$ steht \hfill (Abitur 2025 GK).
\teil Ein Kirchturmdach hat unter anderem die Eckpunkte $B(9 | 0 | 0)$, $C(9 | 9 | 0)$, $G(6 | 9 | 5)$ und $F(9 | 6 | 5)$ (1 LE = 1 m). Die Ebene $L$ enthält $C$, $G$ und $F$; die Ebene $N$ ist parallel zu $L$ und enthält $B$. Weise nach, dass $\vec n = (5 | 5 | 3)$ ein Normalenvektor von $N$ ist, und gib eine Koordinatengleichung von $N$ an \hfill (Abitur 2022 GK).
\teil Gegeben sind $E\colon \vec x = (0 | 1 | 3) + r \cdot (0 | 2 | 1) + s \cdot (1 | 1 | 0)$ und $\vec n = (1 | -1 | 2)$. Weise nach, dass $\vec n$ senkrecht zu $E$ steht \hfill (Abitur 2025 GK).
\teil Gegeben ist $g\colon \vec x = (5 | 2 | -1) + s \cdot (-3 | 0 | 2)$; die Gerade $h$ verläuft parallel zur $y$-Achse und schneidet $g$ im Punkt $(5 | 2 | -1)$. Untersuche, ob $g$ und $h$ senkrecht zueinander verlaufen \hfill (Abitur 2025 GK).
\teil Gegeben ist $g\colon \vec x = (1 | 4 | 2) + t \cdot (2 | -3 | 0)$; die Gerade $h$ verläuft parallel zur $x$-Achse durch $(1 | 4 | 2)$. Untersuche, ob $g$ und $h$ senkrecht zueinander verlaufen \hfill (Abitur 2025 GK).
\teil Gegeben sind $E\colon 4x - 3y + z = 2$ und $F\colon x + 2y + kz = 5$. Bestimme $k$ so, dass $F$ senkrecht zu $E$ steht \hfill (Abitur 2022 GK).
\teil Gegeben sind $E\colon 2y - 5z = 0$ und $F\colon 3x + ay + 2z = 7$. Bestimme $a$ so, dass $F$ senkrecht zu $E$ steht \hfill (Abitur 2022 GK).
\teil Gegeben sind $E\colon x + 2y - z = 3$ und $F\colon 3x - y + 2z = 1$; ihre Schnittgerade hat den Richtungsvektor $(3 | -5 | -7)$. Gib eine Gleichung einer Ebene $H$ an, die senkrecht zu $E$ und zu $F$ steht \hfill (Abitur 2021 GK).
\teil Gegeben sind $E\colon 2x - y + z = 4$ und $F\colon x + y - 3z = 2$; ihre Schnittgerade hat den Richtungsvektor $(2 | 7 | 3)$. Gib eine Gleichung einer Ebene $H$ an, die senkrecht zu $E$ und zu $F$ steht \hfill (Abitur 2021 GK).
\end{teile}
\end{aufgabe}

% Anlauf (ohne Prüfkennung)
\begin{aufgabe}{Flächengleichung des flächenkleinsten Dreiecks begründen – Anlauf}
\begin{teile}
\teil Gegeben sind $R(2 | 0 | 1)$, $S(2 | 0 | 5)$ und die Gerade $t\colon \vec x = (0 | 1 | 3) + \lambda \cdot (1 | 2 | 0)$. Jeder Punkt $T_\lambda$ von $t$ hat von $R$ und $S$ denselben Abstand. Begründe, dass das Dreieck $RST_\lambda$ mit dem kleinsten Flächeninhalt zu dem $\lambda$ gehört, das $(\lambda - 2 | 1 + 2\lambda | 0) \circ (1 | 2 | 0) = 0$ löst, und gib dieses $\lambda$ an.
\end{teile}
\end{aufgabe}

\begin{aufgabe}{Flächengleichung des flächenkleinsten Dreiecks begründen – Prüfungsaufgaben}
\begin{teile}
\teil Gegeben sind $R(0 | 1 | 3)$, $S(0 | 5 | 3)$ und die Gerade $t\colon \vec x = (1 | 3 | 0) + \lambda \cdot (2 | 0 | 1)$. Jeder Punkt $T_\lambda$ von $t$ hat von $R$ und $S$ denselben Abstand. Unter den Dreiecken $RST_\lambda$ hat eines den kleinsten Flächeninhalt. Begründe, dass der zugehörige Wert von $\lambda$ die Gleichung $(1 + 2\lambda | 0 | \lambda - 3) \circ (2 | 0 | 1) = 0$ löst \hfill (Abitur 2020 GK).
\end{teile}
\end{aufgabe}

\clearpage
% Flächeninhalt und Volumen im Raum – bank/flaecheninhalt-und-volumen-im-raum/ (zusammenbau v0.4, Heft)
\einheitenkopf[t10]{Flächeninhalt und Volumen im Raum}
\setcounter{aufgabe}{88}

% Anlauf (ohne Prüfkennung)
\begin{aufgabe}{Dreiecksflächen – Anlauf}
\begin{teile}
\teil Im Dreieck $PQR$ ist $M$ der Mittelpunkt von $PQ$, und $MR$ steht senkrecht auf $PQ$. Was ist $MR$? \\ \kreuz{Höhe zu PQ} \\ \kreuz{schräge Seite}
\teil $P(1 | 1 | 0)$, $Q(4 | 5 | 0)$, $R(-4 | 11 | 0)$, rechtwinklig bei $Q$. Flächeninhalt des Dreiecks $PQR$? A = \leerfeld FE
\teil Die Ebene $3x - 2y + z = 12$ schneidet die $x$-Achse und die $y$-Achse. Diese beiden Schnittpunkte und der Ursprung bilden ein Dreieck. Flächeninhalt? A = \leerfeld FE
\end{teile}
\end{aufgabe}

\begin{aufgabe}{Dreiecksflächen – Prüfungsaufgaben}
\begin{teile}
\teil Der Ursprung $O$, $Q(2 | 4 | 4)$ und $R(5 | 4 | -2)$ bilden ein gleichschenkliges Dreieck mit der Basis $OQ$. Flächeninhalt? \hfill (Abitur 2025 GK) \\ A = \leerfeld FE
\teil Der Ursprung $O$, $P(6 | 8 | 0)$ und $R(3 | 4 | z)$ mit $z > 0$ bilden ein gleichschenkliges Dreieck mit der Basis $OP$ und dem Flächeninhalt 40. Wert von $z$? \hfill (Abitur 2022 GK) \\ z = \leerfeld
\teil $P(1 | 0 | 0)$, $Q(3 | 1 | 0)$, $R(2 | 4 | 1)$; für $T$ gilt $\overrightarrow{OT} = \overrightarrow{OR} - 3 \cdot \overrightarrow{PQ}$. Bestimme $T$ und das Verhältnis der Flächeninhalte von Dreieck $PQR$ und Viereck $PQRT$. \hfill (Abitur 2024 GK)
\end{teile}
\end{aufgabe}

% Anlauf (ohne Prüfkennung)
\begin{aufgabe}{Vierecksflächen – Anlauf}
\begin{teile}
\teil Ein Viereck hat vier gleich lange Seiten, seine Diagonalen sind verschieden lang. Welche Formel für den Flächeninhalt? \\ \kreuz{Seite mal Seite} \\ \kreuz{halbe Summe der parallelen Seiten mal Höhe} \\ \kreuz{halbes Diagonalenprodukt}
\teil $P(1 | 0 | 0)$, $Q(1 | 4 | 3)$, $R(1 | 10 | -5)$. Zeige, dass bei $Q$ ein rechter Winkel liegt, und berechne den Flächeninhalt des Rechtecks $PQRS$. A = \leerfeld FE
\teil Das Viereck mit $P(3 | 4 | 4)$, $Q(-3 | 7 | 4)$, $R(-1 | 2 | 0)$, $S(5 | -1 | 0)$ ist eine Raute. Flächeninhalt? A = \leerfeld FE
\end{teile}
\end{aufgabe}

\begin{aufgabe}{Vierecksflächen – Prüfungsaufgaben}
\begin{teile}
\teil Das Quadrat $PQRS$ mit $P(2 | 1 | 2)$, $Q(8 | 9 | 2)$, $R(0 | 15 | 2)$ und $S(-6 | 7 | 2)$ liegt in der Ebene $z = 2$. Weise nach, dass es den Flächeninhalt 100 hat. \hfill (Abitur 2017 GK)
\teil Das Dach eines Pavillons besteht aus vier gleichen Rauten; eine hat die Ecken $P(6 | 0 | 3)$, $Q(4 | 4 | 5)$, $R(0 | 6 | 3)$, $S(2 | 2 | 1)$ (1 LE = 1 m). Wie groß ist die gesamte Dachfläche? \hfill (Abitur 2022 GK) \\ \leerfeld[m²]
\teil Das Trapez $PQRS$ mit $P(0 | 1 | 1)$, $Q(8 | 1 | 1)$, $R(6 | 4 | 5)$, $S(2 | 4 | 5)$ hat die parallelen Seiten $PQ$ und $SR$ und gleich lange Schenkel. Flächeninhalt? \hfill (Abitur 2019 GK) \\ A = \leerfeld FE
\teil Ein Quadrat liegt in der $x$-$y$-Ebene, eine Ecke ist der Ursprung $O$. Sein Diagonalenschnittpunkt $M$ liegt auf der Geraden $\vec x = (2 | 4 | 4) + t \cdot (2 | 2 | -2)$. Bestimme $M$ und den Flächeninhalt des Quadrats. \hfill (Abitur 2024 GK)
\teil Das Quadrat $PQRS$ mit $P(0 | 1 | 0)$, $Q(8 | 1 | 0)$, $R(8 | 9 | 0)$, $S(0 | 9 | 0)$ wird verändert: $P$ und $S$ werden auf der $y$-Achse gleich weit aufeinander zu verschoben, sodass ein gleichschenkliges Trapez $P'QRS'$ mit drei Vierteln der Quadratfläche entsteht. Bestimme $P'$ und $S'$. \hfill (Abitur 2021 GK)
\teil Ein quadratisches Beet $PQRS$ mit $P(0 | 2 | 0)$, $Q(10 | 2 | 0)$, $R(10 | 12 | 0)$, $S(0 | 12 | 0)$ (1 LE = 1 m) wird verkleinert: $P$ und $S$ rücken auf der $y$-Achse gleich weit aufeinander zu, sodass ein gleichschenkliges Trapez $P'QRS'$ mit drei Fünfteln der Quadratfläche entsteht. Bestimme $P'$ und $S'$. \hfill (Abitur 2021 GK)
\teil Ein Turm hat die quadratische Grundfläche $PQRS$ mit $P(1 | 1 | 0)$, $Q(7 | 1 | 0)$, $R(7 | 7 | 0)$, $S(1 | 7 | 0)$; das Quadrat $KLTU$ mit $K(4 | 1 | 6)$, $L(7 | 4 | 6)$, $T(4 | 7 | 6)$, $U(1 | 4 | 6)$ liegt über den Seitenmitten. Weise nach, dass $PQRS$ doppelt so groß ist wie $KLTU$. \hfill (Abitur 2022 GK)
\teil Ein Sockel hat die quadratische Grundfläche $PQRS$ mit $P(1 | -3 | 0)$, $Q(9 | -3 | 0)$, $R(9 | 5 | 0)$, $S(1 | 5 | 0)$; die Deckfläche $KLTU$ mit $K(5 | -3 | 3)$, $L(9 | 1 | 3)$, $T(5 | 5 | 3)$, $U(1 | 1 | 3)$ liegt über den Seitenmitten. Weise nach, dass $PQRS$ doppelt so groß ist wie $KLTU$. \hfill (Abitur 2022 GK)
\end{teile}
\end{aufgabe}

% Anlauf (ohne Prüfkennung)
\begin{aufgabe}{Ebene Figur: Flächeninhalt des Schattens eines geneigten Rechtecks bei senkrechtem Lichteinfall über eine Querschnittsskizze begründen – Anlauf}
\begin{teile}
\teil Ein rechteckiges Glasdach eines Carports mit der waagerechten Seite 1{,}6\,m und der geneigten Seite 1\,m ist um $40^\circ$ gegen die Waagerechte geneigt; Licht fällt senkrecht auf das Glas. Begründe mit einer Querschnittsskizze, wie lang der Schatten der geneigten Seite auf dem Boden ist, und berechne den Flächeninhalt des Schattens.
\end{teile}
\end{aufgabe}

\begin{aufgabe}{Ebene Figur: Flächeninhalt des Schattens eines geneigten Rechtecks bei senkrechtem Lichteinfall über eine Querschnittsskizze begründen – Prüfungsaufgaben}
\begin{teile}
\teil Ein rechteckiges Vordach ist gegen die Waagerechte geneigt, eine seiner Seiten verläuft waagerecht. Sonnenlicht fällt als parallele Strahlen senkrecht auf die Dachfläche und wirft einen rechteckigen Schatten auf den waagerechten Boden. Begründe mithilfe einer beschrifteten Querschnittsskizze, dass der Schatten einen größeren Flächeninhalt hat als das Vordach. \hfill (Abitur 2017 GK)
\end{teile}
\end{aufgabe}

% Anlauf (ohne Prüfkennung)
\begin{aufgabe}{Volumen und Oberfläche – Anlauf}
\begin{teile}
\teil Ein Zelt über einem quadratischen Boden, alle Dachkanten laufen in einer Spitze zusammen. Volumenformel? \\ \kreuz{mit Faktor ein Drittel} \\ \kreuz{ohne Faktor ein Drittel}
\teil Pyramide mit der Grundfläche $P(1 | 1 | 0)$, $Q(5 | 1 | 0)$, $R(5 | 5 | 0)$, $T(1 | 5 | 0)$ und der Spitze $S(3 | 3 | 6)$. Volumen? V = \leerfeld VE
\teil Eine Pyramide über einem Rechteck mit den Seiten 5 und 8 hat das Volumen 100. Höhe? h = \leerfeld
\end{teile}
\end{aufgabe}

\begin{aufgabe}{Volumen und Oberfläche – Prüfungsaufgaben}
\begin{teile}
\teil Das Dreieck $OPQ$ mit $P(0 | 3 | 4)$ und $Q(6 | 0 | 0)$ ist rechtwinklig in $O$. Die Pyramide $OPQS$ hat das Volumen 40, und es gilt $\overrightarrow{OS} \cdot \overrightarrow{OP} = 0$ und $\overrightarrow{OS} \cdot \overrightarrow{OQ} = 0$. Länge $|OS|$? \hfill (Abitur 2026 GK) \\ |OS| = \leerfeld
\teil Das Quadrat $PQRT$ mit $P(1 | 1 | 2)$, $Q(7 | 1 | 2)$, $R(7 | 7 | 2)$, $T(1 | 7 | 2)$ ist die Grundfläche einer Pyramide $PQRTS$ mit dem Volumen 96. Gib eine mögliche $z$-Koordinate von $S$ an. \hfill (Abitur 2017 GK) \\ z = \leerfeld
\teil Die Pyramide $PQRTS$ hat das Drachenviereck mit $P(0 | 0 | 0)$, $Q(3 | 2 | 0)$, $R(0 | 7 | 0)$, $T(-3 | 2 | 0)$ als Grundfläche und die Spitze $S(0 | 0 | 5)$. Volumen? \hfill (Abitur 2025 GK) \\ V = \leerfeld VE
\teil Ein gerades Prisma hat die Raute mit $P(3 | 0 | 0)$, $Q(0 | 7 | 0)$, $R(-3 | 0 | 0)$, $T(0 | -7 | 0)$ als Grundfläche; die Ecke über $Q$ ist $Q'(0 | 7 | 5)$. Volumen? \hfill (Abitur 2026 GK) \\ V = \leerfeld VE
\teil Ein gerades Prisma hat als Grundfläche ein Dreieck mit den Seiten 5, 12 und 13; seine Mantelfläche ist 150. Höhe des Prismas? \hfill (Abitur 2019 GK) \\ h = \leerfeld
\teil Die Grundfläche der Pyramide $PQRS$ ist ein rechtwinkliges Dreieck $PQR$ mit der Hypotenuse $|PQ| = 13\,\mathrm{cm}$ und der Kathete $|PR| = 5\,\mathrm{cm}$. Die Kante $RS$ steht senkrecht auf der Grundfläche und ist 9\,cm lang. Volumen? \hfill (Abitur 2020 GK) \\ V = \leerfeld[cm³]
\teil Die Grundfläche der Pyramide $PQRS$ ist ein rechtwinkliges Dreieck $PQR$ mit der Hypotenuse $|PQ| = 10\,\mathrm{cm}$ und der Kathete $|PR| = 6\,\mathrm{cm}$. Die Kante $RS$ steht senkrecht auf der Grundfläche und ist 5\,cm lang. Volumen? \hfill (Abitur 2020 GK) \\ V = \leerfeld[cm³]
\end{teile}
\end{aufgabe}

% Anlauf (ohne Prüfkennung)
\begin{aufgabe}{Zusammengesetzt und Verhältnisse – Anlauf}
\begin{teile}
\teil Ein Haus: Quader mit aufgesetztem Satteldach. Wie fasst du sein Volumen? \\ \kreuz{direkt} \\ \kreuz{als Summe} \\ \kreuz{als Differenz}
\teil Ein Quader hat die Grundfläche $PQRU$ mit $P(1 | 1 | 0)$, $Q(7 | 1 | 0)$, $R(7 | 7 | 0)$, $U(1 | 7 | 0)$ und die Höhe 4. Die Ebene durch $Q$, $U$ und $T(1 | 1 | 8)$ zerlegt den Quader. Berechne das Volumen des Teilkörpers, der $P$ enthält, und erläutere den Ansatz. V = \leerfeld VE
\teil Ein Würfel mit der Kantenlänge 8 hat eine Ecke im Ursprung, seine Kanten liegen auf den positiven Achsen. Die Ebene durch die Würfelkante auf der $x$-Achse und den Mittelpunkt $M(8 | 8 | 4)$ einer senkrechten Kante teilt den Würfel. Volumen des kleineren Teilkörpers? V = \leerfeld VE
\end{teile}
\end{aufgabe}

\begin{aufgabe}{Zusammengesetzt und Verhältnisse – Prüfungsaufgaben}
\begin{teile}
\teil Ein Quader hat die Grundfläche $PQRU$ mit $P(2 | 3 | 0)$, $Q(8 | 3 | 0)$, $R(8 | 9 | 0)$, $U(2 | 9 | 0)$ und die Höhe 5. Die Ebene durch $Q$, $U$ und $T(2 | 3 | 10)$ zerlegt den Quader. Berechne das Volumen des Teilkörpers, der $P$ enthält, und erläutere den Ansatz. \hfill (Abitur 2018 GK) \\ V = \leerfeld VE
\teil Ein Sockel hat die Form eines Stumpfs einer geraden quadratischen Pyramide: Grundfläche mit der Seite 6\,m, Deckfläche mit der Seite 3\,m, Höhe 6\,m. Sein Volumen wird mit dem Term $\tfrac{1}{3} \cdot 36 \cdot 12 - \tfrac{1}{3} \cdot 9 \cdot 6$ berechnet. Begründe, dass der Term das Volumen liefert. \hfill (Abitur 2022 GK)
\teil Die Punkte $(9 | 0 | 0)$, $(0 | 9 | 0)$ und $(0 | 0 | 9)$ liegen in der Ebene $x + y + z = 9$; die Punkte $K$, $L$, $N$ liegen je auf einer Koordinatenachse und in der Ebene $3x + 3y + 3z = 12$. Das Volumen des Körpers zwischen den beiden Ebenen im ersten Oktanten wird mit $\tfrac{1}{3} \cdot \tfrac{1}{2} \cdot 9^2 \cdot 9 - \tfrac{1}{3} \cdot \tfrac{1}{2} \cdot k^2 \cdot k$ berechnet. Erläutere den Ansatz und gib $k$ an. \hfill (Abitur 2023 GK)
\teil Ein Quader und eine Pyramide haben dieselbe Grundfläche, die Spitze der Pyramide liegt in der Deckfläche des Quaders. Um wie viel Prozent ist das Volumen des Quaders größer als das der Pyramide? Rechne ohne konkrete Volumina. \hfill (Abitur 2021 GK) \\ \leerfeld \%
\teil Das Dreieck mit den Ecken $(0 | 0)$, $(8 | 0)$ und $(8 | 6)$ rotiert um die $x$-Achse. Oberflächeninhalt des entstehenden Körpers? \hfill (Abitur 2018 GK) \\ O = \leerfeld FE
\teil Eine zylindrische Vorratsdose mit dem Durchmesser 20\,cm fasst 10 Liter. Passt sie in ein Regalfach mit 30\,cm lichter Höhe? \hfill (Abitur 2022 GK)
\end{teile}
\end{aufgabe}

\begin{aufgabe}{Zusammengesetzt und Verhältnisse – Prüfungsaufgaben – weiter}
\begin{teile}
\teil Für $-1 < k < 7$ hat die Pyramide $OP_kQ_kR$ die Ecken $P_k(7 - k | 0 | 0)$, $Q_k(0 | 1 + k | 0)$ und $R(0 | 0 | 3)$. Für welches $k$ ist ihr Volumen am größten? \hfill (Abitur 2023 GK) \\ k = \leerfeld
\teil Für $-4 < k < 8$ hat die Pyramide $OP_kQ_kR$ die Ecken $P_k(4 + k | 0 | 0)$, $Q_k(0 | 8 - k | 0)$ und $R(0 | 0 | 9)$. Für welches $k$ ist ihr Volumen am größten? \hfill (Abitur 2023 GK) \\ k = \leerfeld
\teil Ein Quader hat eine quadratische Grundfläche mit der Seite 6 in der $x$-$y$-Ebene und die Höhe 8. Auf der senkrechten Achse durch seine Mitte liegen $S_k$ in der Höhe $4 - k$ und $Q_k$ in der Höhe $4 + k$ ($0 \le k < 4$). Die Pyramide über der Grundfläche mit der Spitze $S_k$ und die Pyramide über der Deckfläche mit der Spitze $Q_k$ haben zusammen 25\,\% des Quadervolumens. Wert von $k$? \hfill (Abitur 2026 GK) \\ k = \leerfeld
\teil Ein Quader hat eine rechteckige Grundfläche mit den Seiten 4 und 5 in der $x$-$y$-Ebene und die Höhe 10. Auf der senkrechten Achse durch seine Mitte liegen $S_k$ in der Höhe $5 - k$ und $Q_k$ in der Höhe $5 + k$ ($0 \le k < 5$). Die Pyramide über der Grundfläche mit der Spitze $S_k$ und die Pyramide über der Deckfläche mit der Spitze $Q_k$ haben zusammen 20\,\% des Quadervolumens. Wert von $k$? \hfill (Abitur 2026 GK) \\ k = \leerfeld
\teil Aus einem Quader mit der Grundfläche 12 und der Höhe 6 wird die Pyramide über der Grundfläche mit der Spitze $S_t$ in der Höhe $6 - 2t$ herausgenommen ($0 \le t \le 3$). Die Abbildung zeigt zwei Graphen, einer gibt das Volumen des Restkörpers in Abhängigkeit von $t$ an. Welcher? Begründe. \hfill (Abitur 2023 GK) \\ \kreuz{Graph I} \\ \kreuz{Graph II}

\begin{ksys}[xmin=0,xmax=4,ymin=0,ymax=80,xstep=1,ystep=10,xlabel=t,ylabel=V] \funktionab{48+8*\x}{II}{0}{3} \funktionab{48+8/3*\x^2}{I}{0}{3} \end{ksys}
\teil Aus einem Quader mit der Grundfläche 12 und der Höhe 9 wird die Pyramide über der Grundfläche mit der Spitze $S_t$ in der Höhe $9 - 3t$ herausgenommen ($0 \le t \le 3$). Die Abbildung zeigt zwei Graphen, einer gibt das Volumen des Restkörpers in Abhängigkeit von $t$ an. Welcher? Begründe. \hfill (Abitur 2023 GK) \\ \kreuz{Graph I} \\ \kreuz{Graph II}

\begin{ksys}[xmin=0,xmax=4,ymin=0,ymax=120,xstep=1,ystep=20,xlabel=t,ylabel=V] \funktionab{72+12*\x}{I}{0}{3} \funktionab{72+4*\x^2}{II}{0}{3} \end{ksys}
\end{teile}
\end{aufgabe}

\end{document}
